\chapter*{\center \Large ANEXOS}
\addcontentsline{toc}{section}{\bfseries ANEXOS}
\markboth{ANEXOS}{ANEXOS}

\section*{Anexo A. Tablas completas por episodio}

Las tablas siguientes listan, para cada escenario, los parámetros de
entrada y las métricas de salida de todos los episodios de
entrenamiento. Se generan con el script \texttt{tablas\_tex.py}
directamente desde los archivos de resultados por episodio
(\texttt{resultados\_}\allowbreak\texttt{<escenario>.csv}), de modo que
cada número es reproducible con
los comandos del plan de pruebas. Q ini.\ es el número de estados
distintos que la tabla Q ya conocía al empezar el episodio; $r$/dec.\ es
la recompensa media por decisión, medida cada 5~s con el mismo reloj en
todos los controles.

% Generado por tablas_tex.py a partir de los CSV del proyecto. No editar a mano.
\begin{table}[H]
\setstretch{1}
\centering
\scriptsize
\setlength{\tabcolsep}{3pt}
\begin{tabular}{rrrrrrrrrrrr}
\toprule
\multicolumn{4}{c}{Entradas del episodio} & \multicolumn{6}{c}{Salidas del episodio (entrenamiento, $\varepsilon>0$)} & \multicolumn{2}{c}{Política congelada} \\
\cmidrule(lr){1-4}\cmidrule(lr){5-10}\cmidrule(lr){11-12}
Ep. & $\varepsilon$ & Semilla & Q ini. & Retraso & Espera & Cola & $r$/dec. & Cambios & Paradas & Retraso & Espera \\
 &  &  & (estados) & (s/veh) & (s/veh) & (veh) &  & de fase & (/veh) & (s/veh) & (s/veh) \\
\midrule
fijo & 0 & 42 & -- & 22.83 & 9.30 & 5.16 & $-$5.58 & 131 & 0.67 & -- & -- \\
\midrule
1 & 1.000 & 43 & 0 & 80.14 & 30.84 & 14.82 & $-$15.98 & 227 & 3.16 & -- & -- \\
2 & 0.900 & 44 & 81 & 51.59 & 21.21 & 10.90 & $-$11.64 & 219 & 2.16 & -- & -- \\
3 & 0.810 & 45 & 104 & 36.22 & 15.46 & 7.97 & $-$8.51 & 215 & 1.43 & -- & -- \\
4 & 0.729 & 46 & 114 & 33.46 & 14.07 & 7.51 & $-$8.00 & 209 & 1.29 & -- & -- \\
5 & 0.656 & 47 & 119 & 31.21 & 12.33 & 6.62 & $-$7.00 & 222 & 1.27 & 22.15 & 7.26 \\
6 & 0.590 & 48 & 120 & 35.50 & 13.32 & 7.19 & $-$7.61 & 223 & 1.46 & -- & -- \\
7 & 0.531 & 49 & 120 & 24.22 & 8.92 & 4.96 & $-$5.25 & 220 & 0.92 & -- & -- \\
8 & 0.478 & 50 & 123 & 29.17 & 11.04 & 6.03 & $-$6.35 & 219 & 1.17 & -- & -- \\
9 & 0.430 & 51 & 123 & 33.91 & 13.05 & 6.91 & $-$7.29 & 225 & 1.41 & -- & -- \\
10 & 0.387 & 52 & 123 & 27.00 & 9.95 & 5.45 & $-$5.72 & 222 & 1.09 & 21.43 & 7.43 \\
11 & 0.349 & 53 & 123 & 24.91 & 9.03 & 5.04 & $-$5.30 & 215 & 0.98 & -- & -- \\
12 & 0.314 & 54 & 124 & 22.93 & 8.10 & 4.56 & $-$4.74 & 230 & 0.92 & -- & -- \\
13 & 0.282 & 55 & 124 & 21.90 & 7.37 & 4.20 & $-$4.39 & 231 & 0.85 & -- & -- \\
14 & 0.254 & 56 & 124 & 22.11 & 7.51 & 4.24 & $-$4.44 & 222 & 0.86 & -- & -- \\
15 & 0.229 & 57 & 124 & 22.18 & 7.45 & 4.23 & $-$4.41 & 227 & 0.89 & 20.13 & 6.28 \\
16 & 0.206 & 58 & 124 & 20.03 & 6.40 & 3.71 & $-$3.88 & 216 & 0.74 & -- & -- \\
17 & 0.185 & 59 & 125 & 22.28 & 7.77 & 4.40 & $-$4.63 & 215 & 0.88 & -- & -- \\
18 & 0.167 & 60 & 125 & 21.78 & 7.36 & 4.20 & $-$4.39 & 223 & 0.85 & -- & -- \\
19 & 0.150 & 61 & 125 & 26.04 & 9.15 & 5.08 & $-$5.31 & 222 & 1.07 & -- & -- \\
20 & 0.135 & 62 & 126 & 22.10 & 7.59 & 4.31 & $-$4.48 & 229 & 0.88 & 39.18 & 12.66 \\
21 & 0.122 & 63 & 126 & 19.60 & 6.07 & 3.55 & $-$3.69 & 223 & 0.74 & -- & -- \\
22 & 0.109 & 64 & 126 & 21.73 & 7.21 & 4.11 & $-$4.26 & 225 & 0.86 & -- & -- \\
23 & 0.098 & 65 & 127 & 20.76 & 6.71 & 3.86 & $-$3.99 & 225 & 0.80 & -- & -- \\
24 & 0.089 & 66 & 127 & 22.09 & 7.42 & 4.21 & $-$4.37 & 219 & 0.85 & -- & -- \\
25 & 0.080 & 67 & 127 & 19.42 & 6.24 & 3.62 & $-$3.75 & 213 & 0.72 & 20.50 & 6.56 \\
26 & 0.072 & 68 & 127 & 27.84 & 9.91 & 5.45 & $-$5.69 & 230 & 1.18 & -- & -- \\
27 & 0.065 & 69 & 127 & 24.92 & 8.57 & 4.79 & $-$5.03 & 223 & 1.00 & -- & -- \\
28 & 0.058 & 70 & 128 & 21.84 & 7.22 & 4.12 & $-$4.29 & 237 & 0.86 & -- & -- \\
29 & 0.052 & 71 & 128 & 25.07 & 8.80 & 4.92 & $-$5.15 & 229 & 1.03 & -- & -- \\
30 & 0.050 & 72 & 129 & 17.80 & 5.31 & 3.16 & $-$3.30 & 215 & 0.64 & 19.62 & 6.18 \\
31 & 0.050 & 73 & 129 & 20.97 & 6.86 & 3.93 & $-$4.12 & 226 & 0.81 & -- & -- \\
32 & 0.050 & 74 & 129 & 20.66 & 6.73 & 3.83 & $-$3.98 & 229 & 0.81 & -- & -- \\
33 & 0.050 & 75 & 129 & 20.62 & 6.73 & 3.87 & $-$4.02 & 223 & 0.82 & -- & -- \\
34 & 0.050 & 76 & 130 & 19.42 & 6.02 & 3.52 & $-$3.67 & 223 & 0.74 & -- & -- \\
35 & 0.050 & 77 & 130 & 19.66 & 6.23 & 3.62 & $-$3.76 & 219 & 0.74 & 20.02 & 6.36 \\
36 & 0.050 & 78 & 130 & 24.66 & 8.57 & 4.80 & $-$5.02 & 227 & 1.03 & -- & -- \\
37 & 0.050 & 79 & 130 & 19.62 & 6.17 & 3.60 & $-$3.73 & 217 & 0.74 & -- & -- \\
38 & 0.050 & 80 & 130 & 23.06 & 7.92 & 4.47 & $-$4.64 & 225 & 0.94 & -- & -- \\
39 & 0.050 & 81 & 130 & 19.87 & 6.33 & 3.67 & $-$3.81 & 229 & 0.78 & -- & -- \\
40 & 0.050 & 82 & 130 & 21.65 & 7.13 & 4.07 & $-$4.23 & 223 & 0.84 & 21.39 & 7.05 \\
\bottomrule
\end{tabular}
\caption[Escenario cruce: episodios de entrenamiento]{Escenario cruce: parámetros de entrada y métricas de salida por episodio. 40 episodios de entrenamiento con semilla base 42; la fila fijo es el tiempo fijo con la misma semilla base. La política congelada se evalúa cada 5 episodios con una semilla fija ajena al entrenamiento y a la evaluación final.}
\label{tab:cruce-episodios}
\end{table}


% Generado por tablas_tex.py a partir de los CSV del proyecto. No editar a mano.
\begin{table}[H]
\setstretch{1}
\centering
\scriptsize
\setlength{\tabcolsep}{3pt}
\begin{tabular}{rrrrrrrrrrrr}
\toprule
\multicolumn{4}{c}{Entradas del episodio} & \multicolumn{6}{c}{Salidas del episodio (entrenamiento, $\varepsilon>0$)} & \multicolumn{2}{c}{Política congelada} \\
\cmidrule(lr){1-4}\cmidrule(lr){5-10}\cmidrule(lr){11-12}
Ep. & $\varepsilon$ & Semilla & Q ini. & Retraso & Espera & Cola & $r$/dec. & Cambios & Paradas & Retraso & Espera \\
 &  &  & (estados) & (s/veh) & (s/veh) & (veh) &  & de fase & (/veh) & (s/veh) & (s/veh) \\
\midrule
fijo & 0 & 42 & -- & 27.94 & 15.96 & 7.69 & $-$8.36 & 230 & 0.89 & -- & -- \\
\midrule
1 & 1.000 & 43 & 0 & 207.74 & 132.44 & 57.32 & $-$67.25 & 280 & 6.91 & -- & -- \\
2 & 0.960 & 44 & 106 & 207.37 & 134.95 & 57.69 & $-$67.97 & 296 & 6.30 & -- & -- \\
3 & 0.922 & 45 & 128 & 154.32 & 99.04 & 47.59 & $-$56.27 & 256 & 4.74 & -- & -- \\
4 & 0.885 & 46 & 157 & 103.72 & 67.25 & 33.75 & $-$39.39 & 248 & 2.77 & -- & -- \\
5 & 0.849 & 47 & 199 & 89.60 & 59.52 & 28.76 & $-$34.10 & 245 & 2.27 & 209.53 & 153.51 \\
6 & 0.815 & 48 & 241 & 111.92 & 70.79 & 32.27 & $-$37.34 & 284 & 3.66 & -- & -- \\
7 & 0.783 & 49 & 258 & 118.04 & 73.22 & 34.30 & $-$39.97 & 260 & 3.52 & -- & -- \\
8 & 0.751 & 50 & 270 & 59.94 & 39.21 & 18.58 & $-$21.49 & 261 & 1.66 & -- & -- \\
9 & 0.721 & 51 & 274 & 69.69 & 46.06 & 22.02 & $-$25.66 & 255 & 1.82 & -- & -- \\
10 & 0.693 & 52 & 282 & 106.36 & 67.16 & 31.43 & $-$36.27 & 280 & 3.01 & 48.71 & 30.97 \\
11 & 0.665 & 53 & 297 & 52.21 & 34.99 & 17.37 & $-$20.18 & 240 & 1.39 & -- & -- \\
12 & 0.638 & 54 & 307 & 83.47 & 53.19 & 25.56 & $-$29.38 & 267 & 2.41 & -- & -- \\
13 & 0.613 & 55 & 313 & 42.01 & 26.84 & 13.35 & $-$15.28 & 231 & 1.18 & -- & -- \\
14 & 0.588 & 56 & 320 & 47.09 & 30.37 & 15.07 & $-$17.20 & 240 & 1.31 & -- & -- \\
15 & 0.565 & 57 & 328 & 40.80 & 25.90 & 12.70 & $-$14.39 & 256 & 1.19 & 26.65 & 15.05 \\
16 & 0.542 & 58 & 340 & 39.46 & 24.78 & 12.34 & $-$13.88 & 258 & 1.18 & -- & -- \\
17 & 0.520 & 59 & 343 & 39.12 & 24.51 & 12.17 & $-$13.72 & 263 & 1.17 & -- & -- \\
18 & 0.500 & 60 & 348 & 47.08 & 29.93 & 14.80 & $-$16.72 & 263 & 1.37 & -- & -- \\
19 & 0.480 & 61 & 351 & 34.93 & 21.32 & 10.56 & $-$11.79 & 260 & 1.08 & -- & -- \\
20 & 0.460 & 62 & 355 & 38.80 & 23.88 & 11.77 & $-$13.10 & 271 & 1.20 & 33.96 & 20.13 \\
21 & 0.442 & 63 & 359 & 38.99 & 24.10 & 11.61 & $-$12.96 & 267 & 1.18 & -- & -- \\
22 & 0.424 & 64 & 361 & 31.30 & 18.44 & 8.99 & $-$9.98 & 249 & 0.98 & -- & -- \\
23 & 0.407 & 65 & 362 & 34.65 & 20.72 & 10.07 & $-$11.11 & 272 & 1.11 & -- & -- \\
24 & 0.391 & 66 & 362 & 32.28 & 19.13 & 9.39 & $-$10.35 & 259 & 1.04 & -- & -- \\
25 & 0.375 & 67 & 362 & 32.77 & 19.28 & 9.47 & $-$10.39 & 265 & 1.06 & 32.32 & 18.48 \\
26 & 0.360 & 68 & 363 & 39.31 & 24.12 & 12.02 & $-$13.35 & 281 & 1.24 & -- & -- \\
27 & 0.346 & 69 & 367 & 31.48 & 18.47 & 8.93 & $-$9.85 & 266 & 1.01 & -- & -- \\
28 & 0.332 & 70 & 368 & 29.16 & 16.68 & 8.09 & $-$8.89 & 257 & 0.97 & -- & -- \\
29 & 0.319 & 71 & 369 & 33.72 & 20.13 & 9.83 & $-$10.90 & 271 & 1.08 & -- & -- \\
30 & 0.306 & 72 & 372 & 31.39 & 18.35 & 9.02 & $-$9.90 & 266 & 1.03 & 37.14 & 21.75 \\
31 & 0.294 & 73 & 374 & 32.10 & 19.04 & 9.41 & $-$10.38 & 264 & 1.03 & -- & -- \\
32 & 0.282 & 74 & 374 & 30.87 & 17.91 & 8.72 & $-$9.59 & 272 & 1.01 & -- & -- \\
33 & 0.271 & 75 & 374 & 30.37 & 17.52 & 8.58 & $-$9.41 & 261 & 0.99 & -- & -- \\
34 & 0.260 & 76 & 374 & 29.81 & 17.27 & 8.49 & $-$9.35 & 248 & 0.96 & -- & -- \\
35 & 0.250 & 77 & 375 & 28.95 & 16.68 & 8.09 & $-$8.90 & 244 & 0.93 & 26.18 & 14.52 \\
36 & 0.240 & 78 & 375 & 29.90 & 17.22 & 8.47 & $-$9.24 & 267 & 0.99 & -- & -- \\
37 & 0.230 & 79 & 375 & 28.73 & 16.33 & 8.05 & $-$8.78 & 259 & 0.97 & -- & -- \\
38 & 0.221 & 80 & 376 & 33.45 & 19.83 & 9.48 & $-$10.41 & 279 & 1.09 & -- & -- \\
39 & 0.212 & 81 & 376 & 31.79 & 18.66 & 9.25 & $-$10.19 & 267 & 1.05 & -- & -- \\
40 & 0.204 & 82 & 378 & 27.71 & 15.55 & 7.56 & $-$8.22 & 255 & 0.93 & 26.06 & 14.63 \\
\bottomrule
\end{tabular}
\caption[Escenario cruce2: episodios de entrenamiento (parte 1 de 3)]{Escenario cruce2: parámetros de entrada y métricas de salida por episodio (parte 1 de 3). 120 episodios de entrenamiento con semilla base 42; la fila fijo es el tiempo fijo con la misma semilla base. La política congelada se evalúa cada 5 episodios con una semilla fija ajena al entrenamiento y a la evaluación final.}
\label{tab:cruce2-episodios}
\end{table}

\begin{table}[H]
\setstretch{1}
\centering
\scriptsize
\setlength{\tabcolsep}{3pt}
\begin{tabular}{rrrrrrrrrrrr}
\toprule
\multicolumn{4}{c}{Entradas del episodio} & \multicolumn{6}{c}{Salidas del episodio (entrenamiento, $\varepsilon>0$)} & \multicolumn{2}{c}{Política congelada} \\
\cmidrule(lr){1-4}\cmidrule(lr){5-10}\cmidrule(lr){11-12}
Ep. & $\varepsilon$ & Semilla & Q ini. & Retraso & Espera & Cola & $r$/dec. & Cambios & Paradas & Retraso & Espera \\
 &  &  & (estados) & (s/veh) & (s/veh) & (veh) &  & de fase & (/veh) & (s/veh) & (s/veh) \\
\midrule
41 & 0.195 & 83 & 378 & 26.85 & 15.01 & 7.37 & $-$7.99 & 248 & 0.89 & -- & -- \\
42 & 0.188 & 84 & 378 & 27.79 & 15.62 & 7.64 & $-$8.31 & 255 & 0.92 & -- & -- \\
43 & 0.180 & 85 & 378 & 30.56 & 17.63 & 8.59 & $-$9.43 & 264 & 1.00 & -- & -- \\
44 & 0.173 & 86 & 381 & 32.95 & 19.21 & 9.53 & $-$10.44 & 266 & 1.08 & -- & -- \\
45 & 0.166 & 87 & 382 & 31.11 & 18.07 & 8.84 & $-$9.68 & 271 & 1.01 & 26.33 & 14.88 \\
46 & 0.159 & 88 & 383 & 27.10 & 15.12 & 7.31 & $-$7.99 & 255 & 0.91 & -- & -- \\
47 & 0.153 & 89 & 383 & 28.02 & 15.86 & 7.57 & $-$8.27 & 247 & 0.92 & -- & -- \\
48 & 0.147 & 90 & 383 & 27.67 & 15.61 & 7.58 & $-$8.24 & 254 & 0.92 & -- & -- \\
49 & 0.141 & 91 & 383 & 28.56 & 16.27 & 7.98 & $-$8.73 & 254 & 0.94 & -- & -- \\
50 & 0.135 & 92 & 383 & 27.64 & 15.66 & 7.59 & $-$8.32 & 248 & 0.90 & 34.65 & 20.20 \\
51 & 0.130 & 93 & 383 & 26.60 & 14.81 & 7.20 & $-$7.82 & 247 & 0.88 & -- & -- \\
52 & 0.125 & 94 & 383 & 28.53 & 16.24 & 7.99 & $-$8.73 & 246 & 0.93 & -- & -- \\
53 & 0.120 & 95 & 383 & 28.74 & 16.29 & 8.04 & $-$8.77 & 252 & 0.95 & -- & -- \\
54 & 0.115 & 96 & 384 & 26.04 & 14.45 & 7.12 & $-$7.73 & 239 & 0.85 & -- & -- \\
55 & 0.110 & 97 & 384 & 27.12 & 15.23 & 7.46 & $-$8.11 & 248 & 0.91 & 32.38 & 18.73 \\
56 & 0.106 & 98 & 384 & 26.14 & 14.46 & 7.02 & $-$7.59 & 244 & 0.87 & -- & -- \\
57 & 0.102 & 99 & 385 & 26.34 & 14.69 & 7.17 & $-$7.77 & 241 & 0.88 & -- & -- \\
58 & 0.098 & 100 & 385 & 28.50 & 16.01 & 7.85 & $-$8.53 & 267 & 0.95 & -- & -- \\
59 & 0.094 & 101 & 385 & 28.46 & 16.18 & 7.93 & $-$8.69 & 251 & 0.93 & -- & -- \\
60 & 0.090 & 102 & 385 & 28.06 & 15.84 & 7.75 & $-$8.46 & 262 & 0.94 & 47.35 & 28.86 \\
61 & 0.086 & 103 & 385 & 28.28 & 15.98 & 7.89 & $-$8.59 & 259 & 0.94 & -- & -- \\
62 & 0.083 & 104 & 385 & 27.08 & 15.10 & 7.48 & $-$8.19 & 249 & 0.89 & -- & -- \\
63 & 0.080 & 105 & 385 & 25.64 & 14.16 & 6.77 & $-$7.37 & 239 & 0.84 & -- & -- \\
64 & 0.076 & 106 & 385 & 28.41 & 16.18 & 7.84 & $-$8.57 & 255 & 0.93 & -- & -- \\
65 & 0.073 & 107 & 386 & 29.27 & 16.66 & 7.81 & $-$8.55 & 271 & 0.97 & 25.62 & 14.21 \\
66 & 0.070 & 108 & 386 & 27.31 & 15.29 & 7.49 & $-$8.17 & 245 & 0.90 & -- & -- \\
67 & 0.068 & 109 & 386 & 27.22 & 15.15 & 7.43 & $-$8.11 & 245 & 0.90 & -- & -- \\
68 & 0.065 & 110 & 387 & 27.21 & 15.27 & 7.42 & $-$8.06 & 247 & 0.90 & -- & -- \\
69 & 0.062 & 111 & 387 & 27.89 & 15.81 & 7.72 & $-$8.45 & 241 & 0.92 & -- & -- \\
70 & 0.060 & 112 & 387 & 26.78 & 14.92 & 7.26 & $-$7.91 & 235 & 0.89 & 25.90 & 14.28 \\
71 & 0.057 & 113 & 387 & 26.98 & 15.03 & 7.38 & $-$7.99 & 255 & 0.91 & -- & -- \\
72 & 0.055 & 114 & 387 & 28.52 & 16.00 & 7.85 & $-$8.51 & 246 & 0.95 & -- & -- \\
73 & 0.053 & 115 & 387 & 25.51 & 14.08 & 6.86 & $-$7.45 & 226 & 0.84 & -- & -- \\
74 & 0.051 & 116 & 387 & 27.81 & 15.66 & 7.68 & $-$8.41 & 242 & 0.91 & -- & -- \\
75 & 0.050 & 117 & 387 & 24.92 & 13.57 & 6.57 & $-$7.10 & 239 & 0.83 & 68.64 & 45.86 \\
76 & 0.050 & 118 & 387 & 27.85 & 15.62 & 7.69 & $-$8.36 & 257 & 0.92 & -- & -- \\
77 & 0.050 & 119 & 387 & 28.85 & 16.40 & 8.07 & $-$8.81 & 248 & 0.94 & -- & -- \\
78 & 0.050 & 120 & 387 & 27.51 & 15.50 & 7.63 & $-$8.28 & 254 & 0.91 & -- & -- \\
79 & 0.050 & 121 & 387 & 26.71 & 14.90 & 7.34 & $-$8.03 & 239 & 0.87 & -- & -- \\
80 & 0.050 & 122 & 387 & 28.91 & 16.40 & 8.10 & $-$8.80 & 254 & 0.96 & 32.55 & 19.02 \\
\bottomrule
\end{tabular}
\caption[Escenario cruce2: episodios de entrenamiento (parte 2 de 3)]{Escenario cruce2: parámetros de entrada y métricas de salida por episodio (parte 2 de 3). 120 episodios de entrenamiento con semilla base 42; la fila fijo es el tiempo fijo con la misma semilla base. La política congelada se evalúa cada 5 episodios con una semilla fija ajena al entrenamiento y a la evaluación final.}
\label{tab:cruce2-episodios-2}
\end{table}

\begin{table}[H]
\setstretch{1}
\centering
\scriptsize
\setlength{\tabcolsep}{3pt}
\begin{tabular}{rrrrrrrrrrrr}
\toprule
\multicolumn{4}{c}{Entradas del episodio} & \multicolumn{6}{c}{Salidas del episodio (entrenamiento, $\varepsilon>0$)} & \multicolumn{2}{c}{Política congelada} \\
\cmidrule(lr){1-4}\cmidrule(lr){5-10}\cmidrule(lr){11-12}
Ep. & $\varepsilon$ & Semilla & Q ini. & Retraso & Espera & Cola & $r$/dec. & Cambios & Paradas & Retraso & Espera \\
 &  &  & (estados) & (s/veh) & (s/veh) & (veh) &  & de fase & (/veh) & (s/veh) & (s/veh) \\
\midrule
81 & 0.050 & 123 & 387 & 27.44 & 15.32 & 7.45 & $-$8.12 & 253 & 0.90 & -- & -- \\
82 & 0.050 & 124 & 387 & 28.37 & 16.04 & 7.81 & $-$8.52 & 251 & 0.95 & -- & -- \\
83 & 0.050 & 125 & 387 & 26.27 & 14.53 & 7.08 & $-$7.69 & 249 & 0.88 & -- & -- \\
84 & 0.050 & 126 & 387 & 27.55 & 15.58 & 7.61 & $-$8.34 & 233 & 0.90 & -- & -- \\
85 & 0.050 & 127 & 388 & 25.53 & 14.12 & 6.83 & $-$7.43 & 234 & 0.84 & 25.50 & 14.01 \\
86 & 0.050 & 128 & 388 & 25.88 & 14.40 & 6.94 & $-$7.57 & 233 & 0.85 & -- & -- \\
87 & 0.050 & 129 & 388 & 25.26 & 13.82 & 6.77 & $-$7.30 & 238 & 0.85 & -- & -- \\
88 & 0.050 & 130 & 388 & 27.21 & 15.07 & 7.41 & $-$8.02 & 257 & 0.90 & -- & -- \\
89 & 0.050 & 131 & 388 & 28.25 & 16.01 & 7.78 & $-$8.51 & 245 & 0.92 & -- & -- \\
90 & 0.050 & 132 & 389 & 26.26 & 14.46 & 7.03 & $-$7.64 & 253 & 0.88 & 31.05 & 17.48 \\
91 & 0.050 & 133 & 389 & 26.83 & 15.04 & 7.36 & $-$8.04 & 235 & 0.87 & -- & -- \\
92 & 0.050 & 134 & 390 & 25.59 & 14.07 & 6.86 & $-$7.42 & 235 & 0.86 & -- & -- \\
93 & 0.050 & 135 & 390 & 28.51 & 16.16 & 7.97 & $-$8.72 & 247 & 0.93 & -- & -- \\
94 & 0.050 & 136 & 390 & 27.80 & 15.61 & 7.60 & $-$8.26 & 253 & 0.93 & -- & -- \\
95 & 0.050 & 137 & 390 & 29.78 & 17.11 & 8.33 & $-$9.12 & 249 & 0.97 & 25.27 & 13.99 \\
96 & 0.050 & 138 & 390 & 26.39 & 14.77 & 7.17 & $-$7.82 & 236 & 0.86 & -- & -- \\
97 & 0.050 & 139 & 390 & 27.05 & 15.23 & 7.32 & $-$8.01 & 235 & 0.87 & -- & -- \\
98 & 0.050 & 140 & 390 & 26.72 & 14.85 & 7.23 & $-$7.84 & 249 & 0.89 & -- & -- \\
99 & 0.050 & 141 & 390 & 25.49 & 14.02 & 6.89 & $-$7.49 & 233 & 0.83 & -- & -- \\
100 & 0.050 & 142 & 390 & 26.94 & 15.05 & 7.39 & $-$8.08 & 245 & 0.90 & 25.28 & 13.99 \\
101 & 0.050 & 143 & 390 & 25.58 & 14.02 & 6.85 & $-$7.45 & 237 & 0.85 & -- & -- \\
102 & 0.050 & 144 & 390 & 25.46 & 14.03 & 6.83 & $-$7.43 & 233 & 0.85 & -- & -- \\
103 & 0.050 & 145 & 390 & 26.43 & 14.65 & 7.13 & $-$7.81 & 241 & 0.88 & -- & -- \\
104 & 0.050 & 146 & 390 & 27.00 & 15.13 & 7.35 & $-$8.03 & 245 & 0.89 & -- & -- \\
105 & 0.050 & 147 & 390 & 27.42 & 15.49 & 7.54 & $-$8.26 & 237 & 0.89 & 30.19 & 17.22 \\
106 & 0.050 & 148 & 390 & 25.72 & 14.21 & 6.92 & $-$7.52 & 229 & 0.84 & -- & -- \\
107 & 0.050 & 149 & 390 & 27.30 & 15.44 & 7.51 & $-$8.23 & 241 & 0.90 & -- & -- \\
108 & 0.050 & 150 & 390 & 26.21 & 14.57 & 7.13 & $-$7.78 & 233 & 0.86 & -- & -- \\
109 & 0.050 & 151 & 390 & 26.31 & 14.76 & 7.13 & $-$7.81 & 227 & 0.85 & -- & -- \\
110 & 0.050 & 152 & 390 & 25.45 & 14.10 & 6.83 & $-$7.49 & 224 & 0.82 & 29.60 & 16.72 \\
111 & 0.050 & 153 & 390 & 25.64 & 14.16 & 6.89 & $-$7.49 & 236 & 0.85 & -- & -- \\
112 & 0.050 & 154 & 390 & 27.74 & 15.57 & 7.61 & $-$8.29 & 241 & 0.92 & -- & -- \\
113 & 0.050 & 155 & 390 & 25.85 & 14.40 & 6.90 & $-$7.55 & 221 & 0.83 & -- & -- \\
114 & 0.050 & 156 & 390 & 27.46 & 15.46 & 7.48 & $-$8.19 & 233 & 0.89 & -- & -- \\
115 & 0.050 & 157 & 390 & 25.99 & 14.40 & 7.06 & $-$7.66 & 237 & 0.85 & 25.72 & 14.36 \\
116 & 0.050 & 158 & 390 & 25.60 & 14.12 & 6.87 & $-$7.48 & 237 & 0.84 & -- & -- \\
117 & 0.050 & 159 & 390 & 27.24 & 15.25 & 7.44 & $-$8.10 & 245 & 0.91 & -- & -- \\
118 & 0.050 & 160 & 390 & 26.53 & 14.74 & 7.19 & $-$7.80 & 237 & 0.87 & -- & -- \\
119 & 0.050 & 161 & 390 & 28.10 & 15.80 & 7.73 & $-$8.43 & 249 & 0.94 & -- & -- \\
120 & 0.050 & 162 & 390 & 26.53 & 14.81 & 7.19 & $-$7.84 & 238 & 0.87 & 179.00 & 134.62 \\
\bottomrule
\end{tabular}
\caption[Escenario cruce2: episodios de entrenamiento (parte 3 de 3)]{Escenario cruce2: parámetros de entrada y métricas de salida por episodio (parte 3 de 3). 120 episodios de entrenamiento con semilla base 42; la fila fijo es el tiempo fijo con la misma semilla base. La política congelada se evalúa cada 5 episodios con una semilla fija ajena al entrenamiento y a la evaluación final.}
\label{tab:cruce2-episodios-3}
\end{table}


% Generado por tablas_tex.py a partir de los CSV del proyecto. No editar a mano.
\begin{table}[H]
\setstretch{1}
\centering
\scriptsize
\setlength{\tabcolsep}{3pt}
\begin{tabular}{rrrrrrrrrrrr}
\toprule
\multicolumn{4}{c}{Entradas del episodio} & \multicolumn{6}{c}{Salidas del episodio (entrenamiento, $\varepsilon>0$)} & \multicolumn{2}{c}{Política congelada} \\
\cmidrule(lr){1-4}\cmidrule(lr){5-10}\cmidrule(lr){11-12}
Ep. & $\varepsilon$ & Semilla & Q ini. & Retraso & Espera & Cola & $r$/dec. & Cambios & Paradas & Retraso & Espera \\
 &  &  & (estados) & (s/veh) & (s/veh) & (veh) &  & de fase & (/veh) & (s/veh) & (s/veh) \\
\midrule
fijo & 0 & 42 & -- & 31.39 & 14.04 & 7.22 & $-$3.81 & 265 & 1.01 & -- & -- \\
\midrule
1 & 1.000 & 43 & 0 & 86.90 & 38.34 & 17.21 & $-$9.30 & 455 & 3.50 & -- & -- \\
2 & 0.900 & 44 & 133 & 61.99 & 26.70 & 12.51 & $-$6.64 & 450 & 2.60 & -- & -- \\
3 & 0.810 & 45 & 150 & 85.29 & 32.89 & 15.33 & $-$8.12 & 464 & 3.55 & -- & -- \\
4 & 0.729 & 46 & 165 & 54.30 & 21.40 & 10.18 & $-$5.27 & 457 & 2.38 & -- & -- \\
5 & 0.656 & 47 & 172 & 34.42 & 13.08 & 6.77 & $-$3.43 & 447 & 1.52 & 24.74 & 8.19 \\
6 & 0.590 & 48 & 176 & 47.34 & 18.83 & 9.33 & $-$4.86 & 455 & 2.09 & -- & -- \\
7 & 0.531 & 49 & 179 & 42.04 & 16.74 & 8.43 & $-$4.37 & 444 & 1.79 & -- & -- \\
8 & 0.478 & 50 & 183 & 37.35 & 14.26 & 7.28 & $-$3.72 & 450 & 1.68 & -- & -- \\
9 & 0.430 & 51 & 185 & 55.86 & 20.55 & 10.19 & $-$5.27 & 452 & 2.40 & -- & -- \\
10 & 0.387 & 52 & 188 & 40.19 & 14.72 & 7.48 & $-$3.82 & 456 & 1.78 & 30.09 & 10.90 \\
11 & 0.349 & 53 & 192 & 37.30 & 14.05 & 7.29 & $-$3.71 & 456 & 1.66 & -- & -- \\
12 & 0.314 & 54 & 193 & 26.84 & 9.69 & 5.15 & $-$2.57 & 442 & 1.15 & -- & -- \\
13 & 0.282 & 55 & 193 & 30.70 & 11.35 & 5.91 & $-$2.98 & 461 & 1.36 & -- & -- \\
14 & 0.254 & 56 & 194 & 30.45 & 11.55 & 6.02 & $-$3.06 & 466 & 1.33 & -- & -- \\
15 & 0.229 & 57 & 194 & 35.30 & 13.02 & 6.75 & $-$3.40 & 472 & 1.64 & 30.45 & 10.74 \\
16 & 0.206 & 58 & 195 & 31.78 & 11.81 & 6.18 & $-$3.10 & 450 & 1.39 & -- & -- \\
17 & 0.185 & 59 & 196 & 27.17 & 9.41 & 5.03 & $-$2.50 & 462 & 1.22 & -- & -- \\
18 & 0.167 & 60 & 196 & 29.16 & 10.79 & 5.70 & $-$2.85 & 465 & 1.28 & -- & -- \\
19 & 0.150 & 61 & 197 & 28.32 & 9.93 & 5.27 & $-$2.63 & 462 & 1.25 & -- & -- \\
20 & 0.135 & 62 & 197 & 29.06 & 10.06 & 5.30 & $-$2.64 & 472 & 1.31 & 23.89 & 7.91 \\
21 & 0.122 & 63 & 197 & 30.29 & 10.94 & 5.75 & $-$2.88 & 459 & 1.37 & -- & -- \\
22 & 0.109 & 64 & 197 & 27.97 & 10.87 & 5.75 & $-$2.93 & 420 & 1.16 & -- & -- \\
23 & 0.098 & 65 & 198 & 30.60 & 10.87 & 5.70 & $-$2.85 & 470 & 1.39 & -- & -- \\
24 & 0.089 & 66 & 201 & 37.55 & 13.27 & 6.80 & $-$3.45 & 464 & 1.67 & -- & -- \\
25 & 0.080 & 67 & 201 & 27.95 & 9.90 & 5.28 & $-$2.63 & 441 & 1.22 & 30.14 & 12.26 \\
26 & 0.072 & 68 & 201 & 27.77 & 10.02 & 5.29 & $-$2.64 & 452 & 1.21 & -- & -- \\
27 & 0.065 & 69 & 201 & 26.39 & 9.08 & 4.86 & $-$2.42 & 464 & 1.14 & -- & -- \\
28 & 0.058 & 70 & 201 & 30.55 & 10.55 & 5.56 & $-$2.77 & 462 & 1.42 & -- & -- \\
29 & 0.052 & 71 & 201 & 38.56 & 12.95 & 6.65 & $-$3.38 & 448 & 1.75 & -- & -- \\
30 & 0.050 & 72 & 201 & 42.91 & 14.94 & 7.62 & $-$3.87 & 448 & 1.96 & 27.97 & 9.69 \\
31 & 0.050 & 73 & 203 & 26.17 & 9.07 & 4.88 & $-$2.43 & 444 & 1.12 & -- & -- \\
32 & 0.050 & 74 & 203 & 32.18 & 10.78 & 5.71 & $-$2.84 & 458 & 1.45 & -- & -- \\
33 & 0.050 & 75 & 203 & 26.34 & 9.12 & 4.89 & $-$2.43 & 466 & 1.15 & -- & -- \\
34 & 0.050 & 76 & 203 & 26.81 & 9.13 & 4.90 & $-$2.43 & 463 & 1.18 & -- & -- \\
35 & 0.050 & 77 & 203 & 28.43 & 9.89 & 5.23 & $-$2.61 & 444 & 1.25 & 24.34 & 7.98 \\
36 & 0.050 & 78 & 203 & 26.78 & 9.43 & 5.01 & $-$2.49 & 451 & 1.17 & -- & -- \\
37 & 0.050 & 79 & 203 & 26.24 & 8.86 & 4.76 & $-$2.35 & 447 & 1.16 & -- & -- \\
38 & 0.050 & 80 & 203 & 27.02 & 9.29 & 4.96 & $-$2.45 & 455 & 1.20 & -- & -- \\
39 & 0.050 & 81 & 204 & 29.00 & 10.25 & 5.40 & $-$2.70 & 456 & 1.29 & -- & -- \\
40 & 0.050 & 82 & 204 & 29.29 & 10.60 & 5.59 & $-$2.81 & 448 & 1.27 & 30.14 & 10.81 \\
\bottomrule
\end{tabular}
\caption[Escenario corredor: episodios de entrenamiento]{Escenario corredor: parámetros de entrada y métricas de salida por episodio. 40 episodios de entrenamiento con semilla base 42; la fila fijo es el tiempo fijo con la misma semilla base. La política congelada se evalúa cada 5 episodios con una semilla fija ajena al entrenamiento y a la evaluación final.}
\label{tab:corredor-episodios}
\end{table}


% Generado por tablas_tex.py a partir de los CSV del proyecto. No editar a mano.
\begin{table}[H]
\setstretch{1}
\centering
\scriptsize
\setlength{\tabcolsep}{3pt}
\begin{tabular}{rrrrrrrrrrrr}
\toprule
\multicolumn{4}{c}{Entradas del episodio} & \multicolumn{6}{c}{Salidas del episodio (entrenamiento, $\varepsilon>0$)} & \multicolumn{2}{c}{Política congelada} \\
\cmidrule(lr){1-4}\cmidrule(lr){5-10}\cmidrule(lr){11-12}
Ep. & $\varepsilon$ & Semilla & Q ini. & Retraso & Espera & Cola & $r$/dec. & Cambios & Paradas & Retraso & Espera \\
 &  &  & (estados) & (s/veh) & (s/veh) & (veh) &  & de fase & (/veh) & (s/veh) & (s/veh) \\
\midrule
fijo & 0 & 42 & -- & 22.23 & 7.35 & 4.03 & $-$1.43 & 489 & 0.60 & -- & -- \\
\midrule
1 & 1.000 & 43 & 0 & 30.94 & 12.28 & 6.73 & $-$2.41 & 629 & 1.15 & -- & -- \\
2 & 0.950 & 44 & 88 & 30.71 & 12.15 & 6.70 & $-$2.40 & 649 & 1.16 & -- & -- \\
3 & 0.902 & 45 & 108 & 30.71 & 12.23 & 6.70 & $-$2.40 & 637 & 1.14 & -- & -- \\
4 & 0.857 & 46 & 124 & 29.33 & 11.19 & 6.16 & $-$2.19 & 662 & 1.12 & -- & -- \\
5 & 0.815 & 47 & 128 & 29.32 & 10.97 & 6.14 & $-$2.18 & 660 & 1.13 & 25.47 & 8.07 \\
6 & 0.774 & 48 & 132 & 27.36 & 9.71 & 5.31 & $-$1.87 & 683 & 1.06 & -- & -- \\
7 & 0.735 & 49 & 134 & 27.81 & 9.89 & 5.41 & $-$1.91 & 678 & 1.09 & -- & -- \\
8 & 0.698 & 50 & 136 & 28.97 & 10.89 & 5.04 & $-$1.77 & 684 & 1.12 & -- & -- \\
9 & 0.663 & 51 & 139 & 28.04 & 10.03 & 5.33 & $-$1.87 & 690 & 1.08 & -- & -- \\
10 & 0.630 & 52 & 139 & 27.96 & 10.06 & 5.47 & $-$1.92 & 683 & 1.08 & 25.32 & 8.08 \\
11 & 0.599 & 53 & 140 & 27.50 & 9.61 & 5.44 & $-$1.90 & 688 & 1.09 & -- & -- \\
12 & 0.569 & 54 & 140 & 26.65 & 8.96 & 4.82 & $-$1.68 & 724 & 1.04 & -- & -- \\
13 & 0.540 & 55 & 140 & 26.96 & 9.29 & 5.01 & $-$1.76 & 717 & 1.08 & -- & -- \\
14 & 0.513 & 56 & 140 & 26.11 & 8.76 & 4.74 & $-$1.66 & 706 & 1.03 & -- & -- \\
15 & 0.488 & 57 & 140 & 27.57 & 9.68 & 5.30 & $-$1.87 & 708 & 1.11 & 24.88 & 7.67 \\
16 & 0.463 & 58 & 141 & 27.36 & 9.45 & 5.19 & $-$1.80 & 720 & 1.12 & -- & -- \\
17 & 0.440 & 59 & 141 & 26.94 & 9.22 & 5.01 & $-$1.76 & 713 & 1.10 & -- & -- \\
18 & 0.418 & 60 & 142 & 25.62 & 8.29 & 4.48 & $-$1.57 & 714 & 1.02 & -- & -- \\
19 & 0.397 & 61 & 142 & 26.69 & 8.98 & 4.72 & $-$1.66 & 730 & 1.08 & -- & -- \\
20 & 0.377 & 62 & 142 & 25.58 & 8.38 & 4.47 & $-$1.56 & 737 & 1.04 & 24.78 & 7.93 \\
21 & 0.358 & 63 & 142 & 25.81 & 8.32 & 4.72 & $-$1.65 & 725 & 1.06 & -- & -- \\
22 & 0.341 & 64 & 142 & 24.97 & 7.60 & 4.17 & $-$1.44 & 775 & 1.03 & -- & -- \\
23 & 0.324 & 65 & 143 & 26.29 & 8.77 & 4.89 & $-$1.71 & 748 & 1.09 & -- & -- \\
24 & 0.307 & 66 & 143 & 25.21 & 7.90 & 4.30 & $-$1.49 & 751 & 1.02 & -- & -- \\
25 & 0.292 & 67 & 143 & 26.52 & 8.95 & 4.92 & $-$1.71 & 724 & 1.09 & 24.29 & 7.19 \\
26 & 0.277 & 68 & 144 & 25.63 & 8.22 & 4.51 & $-$1.57 & 728 & 1.03 & -- & -- \\
27 & 0.264 & 69 & 144 & 25.73 & 8.22 & 4.46 & $-$1.54 & 719 & 1.06 & -- & -- \\
28 & 0.250 & 70 & 144 & 25.24 & 8.17 & 4.47 & $-$1.55 & 732 & 1.05 & -- & -- \\
29 & 0.238 & 71 & 144 & 25.54 & 8.32 & 4.31 & $-$1.48 & 748 & 1.05 & -- & -- \\
30 & 0.226 & 72 & 144 & 25.91 & 8.69 & 4.47 & $-$1.54 & 736 & 1.08 & 24.27 & 7.31 \\
31 & 0.215 & 73 & 144 & 25.08 & 7.87 & 4.36 & $-$1.51 & 744 & 1.03 & -- & -- \\
32 & 0.204 & 74 & 144 & 24.21 & 7.60 & 4.15 & $-$1.43 & 737 & 0.97 & -- & -- \\
33 & 0.194 & 75 & 144 & 25.64 & 8.36 & 4.51 & $-$1.56 & 751 & 1.04 & -- & -- \\
34 & 0.184 & 76 & 145 & 25.08 & 8.02 & 4.44 & $-$1.55 & 739 & 1.03 & -- & -- \\
35 & 0.175 & 77 & 145 & 24.72 & 7.72 & 4.26 & $-$1.49 & 707 & 1.00 & 24.44 & 7.65 \\
36 & 0.166 & 78 & 145 & 24.47 & 7.58 & 4.28 & $-$1.49 & 722 & 1.00 & -- & -- \\
37 & 0.158 & 79 & 145 & 25.03 & 7.94 & 4.21 & $-$1.46 & 732 & 1.00 & -- & -- \\
38 & 0.150 & 80 & 145 & 24.36 & 7.39 & 4.24 & $-$1.47 & 742 & 0.97 & -- & -- \\
39 & 0.142 & 81 & 145 & 25.99 & 8.85 & 4.35 & $-$1.50 & 741 & 1.05 & -- & -- \\
40 & 0.135 & 82 & 145 & 24.41 & 7.54 & 4.17 & $-$1.45 & 733 & 0.97 & 23.91 & 7.29 \\
\bottomrule
\end{tabular}
\caption[Escenario red: episodios de entrenamiento (parte 1 de 3)]{Escenario red: parámetros de entrada y métricas de salida por episodio (parte 1 de 3). 100 episodios de entrenamiento con semilla base 42; la fila fijo es el tiempo fijo con la misma semilla base. La política congelada se evalúa cada 5 episodios con una semilla fija ajena al entrenamiento y a la evaluación final.}
\label{tab:red-episodios}
\end{table}

\begin{table}[H]
\setstretch{1}
\centering
\scriptsize
\setlength{\tabcolsep}{3pt}
\begin{tabular}{rrrrrrrrrrrr}
\toprule
\multicolumn{4}{c}{Entradas del episodio} & \multicolumn{6}{c}{Salidas del episodio (entrenamiento, $\varepsilon>0$)} & \multicolumn{2}{c}{Política congelada} \\
\cmidrule(lr){1-4}\cmidrule(lr){5-10}\cmidrule(lr){11-12}
Ep. & $\varepsilon$ & Semilla & Q ini. & Retraso & Espera & Cola & $r$/dec. & Cambios & Paradas & Retraso & Espera \\
 &  &  & (estados) & (s/veh) & (s/veh) & (veh) &  & de fase & (/veh) & (s/veh) & (s/veh) \\
\midrule
41 & 0.129 & 83 & 145 & 25.14 & 8.10 & 4.36 & $-$1.52 & 747 & 1.04 & -- & -- \\
42 & 0.122 & 84 & 145 & 24.93 & 7.91 & 4.33 & $-$1.50 & 766 & 0.99 & -- & -- \\
43 & 0.116 & 85 & 145 & 25.19 & 8.14 & 4.47 & $-$1.55 & 754 & 1.01 & -- & -- \\
44 & 0.110 & 86 & 145 & 24.91 & 7.84 & 4.43 & $-$1.53 & 766 & 1.01 & -- & -- \\
45 & 0.105 & 87 & 145 & 25.28 & 8.08 & 4.27 & $-$1.48 & 769 & 1.06 & 25.49 & 7.90 \\
46 & 0.099 & 88 & 145 & 25.13 & 8.05 & 4.16 & $-$1.45 & 762 & 1.03 & -- & -- \\
47 & 0.094 & 89 & 145 & 24.81 & 7.93 & 4.23 & $-$1.47 & 761 & 1.00 & -- & -- \\
48 & 0.090 & 90 & 145 & 24.59 & 7.54 & 4.27 & $-$1.47 & 768 & 1.00 & -- & -- \\
49 & 0.085 & 91 & 145 & 24.36 & 7.64 & 4.13 & $-$1.43 & 757 & 0.98 & -- & -- \\
50 & 0.081 & 92 & 145 & 24.85 & 7.80 & 4.35 & $-$1.50 & 746 & 1.03 & 24.00 & 7.03 \\
51 & 0.077 & 93 & 145 & 24.08 & 7.34 & 4.03 & $-$1.39 & 756 & 0.96 & -- & -- \\
52 & 0.073 & 94 & 145 & 25.65 & 8.49 & 4.20 & $-$1.46 & 764 & 1.04 & -- & -- \\
53 & 0.069 & 95 & 145 & 24.99 & 8.15 & 4.33 & $-$1.51 & 765 & 1.01 & -- & -- \\
54 & 0.066 & 96 & 145 & 25.27 & 7.93 & 4.36 & $-$1.50 & 753 & 1.04 & -- & -- \\
55 & 0.063 & 97 & 145 & 24.73 & 7.73 & 4.20 & $-$1.45 & 768 & 1.01 & 24.72 & 7.71 \\
56 & 0.060 & 98 & 145 & 25.12 & 8.16 & 4.28 & $-$1.49 & 752 & 1.00 & -- & -- \\
57 & 0.057 & 99 & 145 & 24.60 & 7.87 & 3.98 & $-$1.38 & 748 & 0.99 & -- & -- \\
58 & 0.054 & 100 & 145 & 25.32 & 7.78 & 4.32 & $-$1.50 & 769 & 1.08 & -- & -- \\
59 & 0.051 & 101 & 145 & 26.26 & 8.86 & 4.24 & $-$1.45 & 773 & 1.12 & -- & -- \\
60 & 0.050 & 102 & 145 & 25.00 & 7.89 & 4.38 & $-$1.52 & 757 & 1.03 & 24.45 & 7.69 \\
61 & 0.050 & 103 & 146 & 24.96 & 8.11 & 4.32 & $-$1.50 & 747 & 1.02 & -- & -- \\
62 & 0.050 & 104 & 146 & 23.75 & 7.29 & 4.14 & $-$1.42 & 738 & 0.94 & -- & -- \\
63 & 0.050 & 105 & 146 & 23.99 & 7.74 & 4.03 & $-$1.39 & 735 & 0.91 & -- & -- \\
64 & 0.050 & 106 & 146 & 24.52 & 7.82 & 4.34 & $-$1.52 & 749 & 0.97 & -- & -- \\
65 & 0.050 & 107 & 146 & 25.31 & 8.04 & 4.36 & $-$1.50 & 761 & 1.03 & 24.36 & 7.51 \\
66 & 0.050 & 108 & 146 & 25.10 & 8.07 & 4.32 & $-$1.50 & 747 & 1.01 & -- & -- \\
67 & 0.050 & 109 & 146 & 24.45 & 7.52 & 4.29 & $-$1.49 & 751 & 0.99 & -- & -- \\
68 & 0.050 & 110 & 146 & 24.80 & 7.91 & 4.25 & $-$1.47 & 736 & 1.00 & -- & -- \\
69 & 0.050 & 111 & 146 & 23.89 & 7.29 & 4.03 & $-$1.40 & 749 & 0.94 & -- & -- \\
70 & 0.050 & 112 & 146 & 24.45 & 7.56 & 4.22 & $-$1.46 & 748 & 0.97 & 24.23 & 7.40 \\
71 & 0.050 & 113 & 148 & 25.40 & 8.30 & 4.34 & $-$1.49 & 751 & 1.06 & -- & -- \\
72 & 0.050 & 114 & 148 & 24.79 & 7.68 & 4.24 & $-$1.47 & 770 & 1.01 & -- & -- \\
73 & 0.050 & 115 & 148 & 24.20 & 7.58 & 4.25 & $-$1.47 & 764 & 1.00 & -- & -- \\
74 & 0.050 & 116 & 148 & 24.74 & 7.52 & 4.32 & $-$1.50 & 758 & 1.01 & -- & -- \\
75 & 0.050 & 117 & 148 & 24.86 & 7.74 & 4.25 & $-$1.47 & 750 & 1.02 & 23.85 & 7.15 \\
76 & 0.050 & 118 & 148 & 24.30 & 7.49 & 4.02 & $-$1.40 & 759 & 0.94 & -- & -- \\
77 & 0.050 & 119 & 148 & 24.24 & 7.45 & 4.09 & $-$1.42 & 751 & 0.98 & -- & -- \\
78 & 0.050 & 120 & 148 & 24.98 & 8.09 & 4.16 & $-$1.44 & 756 & 1.00 & -- & -- \\
79 & 0.050 & 121 & 148 & 23.67 & 7.29 & 4.00 & $-$1.39 & 759 & 0.94 & -- & -- \\
80 & 0.050 & 122 & 148 & 24.52 & 7.57 & 4.26 & $-$1.47 & 747 & 0.97 & 23.77 & 7.03 \\
\bottomrule
\end{tabular}
\caption[Escenario red: episodios de entrenamiento (parte 2 de 3)]{Escenario red: parámetros de entrada y métricas de salida por episodio (parte 2 de 3). 100 episodios de entrenamiento con semilla base 42; la fila fijo es el tiempo fijo con la misma semilla base. La política congelada se evalúa cada 5 episodios con una semilla fija ajena al entrenamiento y a la evaluación final.}
\label{tab:red-episodios-2}
\end{table}

\begin{table}[H]
\setstretch{1}
\centering
\scriptsize
\setlength{\tabcolsep}{3pt}
\begin{tabular}{rrrrrrrrrrrr}
\toprule
\multicolumn{4}{c}{Entradas del episodio} & \multicolumn{6}{c}{Salidas del episodio (entrenamiento, $\varepsilon>0$)} & \multicolumn{2}{c}{Política congelada} \\
\cmidrule(lr){1-4}\cmidrule(lr){5-10}\cmidrule(lr){11-12}
Ep. & $\varepsilon$ & Semilla & Q ini. & Retraso & Espera & Cola & $r$/dec. & Cambios & Paradas & Retraso & Espera \\
 &  &  & (estados) & (s/veh) & (s/veh) & (veh) &  & de fase & (/veh) & (s/veh) & (s/veh) \\
\midrule
81 & 0.050 & 123 & 148 & 24.61 & 7.71 & 4.26 & $-$1.47 & 748 & 1.00 & -- & -- \\
82 & 0.050 & 124 & 148 & 24.26 & 7.55 & 4.13 & $-$1.43 & 769 & 0.98 & -- & -- \\
83 & 0.050 & 125 & 148 & 24.80 & 7.83 & 4.39 & $-$1.54 & 772 & 1.02 & -- & -- \\
84 & 0.050 & 126 & 148 & 24.00 & 7.44 & 4.03 & $-$1.39 & 737 & 0.95 & -- & -- \\
85 & 0.050 & 127 & 148 & 24.59 & 7.73 & 4.25 & $-$1.47 & 748 & 1.00 & 24.46 & 7.56 \\
86 & 0.050 & 128 & 148 & 24.89 & 7.97 & 4.26 & $-$1.48 & 740 & 0.97 & -- & -- \\
87 & 0.050 & 129 & 148 & 24.71 & 7.81 & 4.29 & $-$1.49 & 738 & 0.99 & -- & -- \\
88 & 0.050 & 130 & 148 & 23.99 & 7.22 & 4.15 & $-$1.45 & 731 & 0.96 & -- & -- \\
89 & 0.050 & 131 & 149 & 24.05 & 7.33 & 4.09 & $-$1.42 & 758 & 0.95 & -- & -- \\
90 & 0.050 & 132 & 149 & 24.63 & 7.69 & 4.22 & $-$1.47 & 766 & 0.99 & 24.38 & 7.56 \\
91 & 0.050 & 133 & 149 & 24.64 & 7.76 & 4.36 & $-$1.52 & 755 & 0.99 & -- & -- \\
92 & 0.050 & 134 & 149 & 24.70 & 7.62 & 4.24 & $-$1.47 & 748 & 1.03 & -- & -- \\
93 & 0.050 & 135 & 149 & 25.71 & 8.58 & 4.47 & $-$1.53 & 758 & 1.04 & -- & -- \\
94 & 0.050 & 136 & 149 & 24.56 & 7.89 & 4.15 & $-$1.44 & 755 & 0.96 & -- & -- \\
95 & 0.050 & 137 & 149 & 24.30 & 7.65 & 4.11 & $-$1.43 & 758 & 0.95 & 24.85 & 7.73 \\
96 & 0.050 & 138 & 149 & 25.11 & 7.99 & 4.37 & $-$1.52 & 775 & 1.03 & -- & -- \\
97 & 0.050 & 139 & 149 & 24.62 & 7.69 & 4.16 & $-$1.45 & 775 & 0.98 & -- & -- \\
98 & 0.050 & 140 & 149 & 25.48 & 8.22 & 4.15 & $-$1.43 & 774 & 1.02 & -- & -- \\
99 & 0.050 & 141 & 149 & 24.94 & 7.98 & 4.45 & $-$1.54 & 765 & 1.01 & -- & -- \\
100 & 0.050 & 142 & 149 & 25.36 & 8.16 & 4.58 & $-$1.59 & 761 & 1.05 & 24.84 & 7.58 \\
\bottomrule
\end{tabular}
\caption[Escenario red: episodios de entrenamiento (parte 3 de 3)]{Escenario red: parámetros de entrada y métricas de salida por episodio (parte 3 de 3). 100 episodios de entrenamiento con semilla base 42; la fila fijo es el tiempo fijo con la misma semilla base. La política congelada se evalúa cada 5 episodios con una semilla fija ajena al entrenamiento y a la evaluación final.}
\label{tab:red-episodios-3}
\end{table}


% Generado por tablas_tex.py a partir de los CSV del proyecto. No editar a mano.
\begin{table}[H]
\setstretch{1}
\centering
\scriptsize
\setlength{\tabcolsep}{3pt}
\begin{tabular}{rrrrrrrrrrrr}
\toprule
\multicolumn{4}{c}{Entradas del episodio} & \multicolumn{6}{c}{Salidas del episodio (entrenamiento, $\varepsilon>0$)} & \multicolumn{2}{c}{Política congelada} \\
\cmidrule(lr){1-4}\cmidrule(lr){5-10}\cmidrule(lr){11-12}
Ep. & $\varepsilon$ & Semilla & Q ini. & Retraso & Espera & Cola & $r$/dec. & Cambios & Paradas & Retraso & Espera \\
 &  &  & (estados) & (s/veh) & (s/veh) & (veh) &  & de fase & (/veh) & (s/veh) & (s/veh) \\
\midrule
fijo & 0 & 42 & -- & 157.82 & 64.97 & 46.13 & $-$24.76 & 448 & 6.48 & -- & -- \\
\midrule
1 & 1.000 & 43 & 0 & 212.01 & 86.56 & 51.37 & $-$28.02 & 660 & 8.34 & -- & -- \\
2 & 0.900 & 44 & 274 & 207.79 & 80.75 & 48.63 & $-$26.39 & 680 & 8.28 & -- & -- \\
3 & 0.810 & 45 & 364 & 191.94 & 67.82 & 41.38 & $-$22.16 & 678 & 7.86 & -- & -- \\
4 & 0.729 & 46 & 408 & 184.27 & 66.58 & 40.42 & $-$21.63 & 685 & 7.54 & -- & -- \\
5 & 0.656 & 47 & 448 & 182.31 & 65.17 & 39.75 & $-$21.37 & 665 & 7.54 & 155.54 & 60.41 \\
6 & 0.590 & 48 & 474 & 178.22 & 59.52 & 36.85 & $-$19.66 & 654 & 7.21 & -- & -- \\
7 & 0.531 & 49 & 491 & 186.30 & 63.32 & 38.64 & $-$20.59 & 672 & 7.77 & -- & -- \\
8 & 0.478 & 50 & 520 & 163.34 & 55.53 & 34.69 & $-$18.46 & 652 & 6.72 & -- & -- \\
9 & 0.430 & 51 & 530 & 181.00 & 62.44 & 38.11 & $-$20.37 & 686 & 7.50 & -- & -- \\
10 & 0.387 & 52 & 538 & 176.44 & 58.67 & 36.16 & $-$19.30 & 671 & 7.21 & 169.61 & 63.92 \\
11 & 0.349 & 53 & 554 & 176.23 & 60.52 & 37.92 & $-$20.44 & 639 & 6.98 & -- & -- \\
12 & 0.314 & 54 & 562 & 171.65 & 55.01 & 35.20 & $-$18.77 & 616 & 7.02 & -- & -- \\
13 & 0.282 & 55 & 568 & 178.44 & 59.70 & 37.34 & $-$20.12 & 630 & 7.18 & -- & -- \\
14 & 0.254 & 56 & 577 & 162.07 & 51.31 & 31.97 & $-$17.05 & 673 & 6.58 & -- & -- \\
15 & 0.229 & 57 & 585 & 167.37 & 51.47 & 32.18 & $-$17.13 & 656 & 6.73 & 173.82 & 54.56 \\
16 & 0.206 & 58 & 594 & 183.91 & 60.20 & 38.60 & $-$20.66 & 620 & 7.43 & -- & -- \\
17 & 0.185 & 59 & 615 & 169.89 & 55.47 & 35.49 & $-$19.02 & 612 & 6.97 & -- & -- \\
18 & 0.167 & 60 & 621 & 179.64 & 59.86 & 37.91 & $-$20.42 & 622 & 7.18 & -- & -- \\
19 & 0.150 & 61 & 635 & 175.45 & 57.92 & 37.78 & $-$20.10 & 602 & 7.15 & -- & -- \\
20 & 0.135 & 62 & 642 & 167.83 & 54.47 & 34.49 & $-$18.46 & 619 & 6.73 & 156.02 & 60.51 \\
21 & 0.122 & 63 & 649 & 155.72 & 50.63 & 33.60 & $-$17.91 & 594 & 6.43 & -- & -- \\
22 & 0.109 & 64 & 656 & 178.57 & 60.49 & 38.39 & $-$20.62 & 620 & 7.24 & -- & -- \\
23 & 0.098 & 65 & 665 & 185.45 & 64.66 & 40.25 & $-$21.71 & 636 & 7.48 & -- & -- \\
24 & 0.089 & 66 & 667 & 186.79 & 66.08 & 42.31 & $-$22.88 & 580 & 7.37 & -- & -- \\
25 & 0.080 & 67 & 672 & 177.39 & 56.11 & 35.64 & $-$19.06 & 622 & 7.17 & 175.29 & 61.76 \\
26 & 0.072 & 68 & 676 & 175.54 & 58.34 & 36.01 & $-$19.42 & 658 & 7.04 & -- & -- \\
27 & 0.065 & 69 & 677 & 196.62 & 63.58 & 38.80 & $-$20.74 & 672 & 7.94 & -- & -- \\
28 & 0.058 & 70 & 680 & 178.02 & 60.53 & 39.55 & $-$21.33 & 598 & 6.99 & -- & -- \\
29 & 0.052 & 71 & 685 & 183.02 & 64.68 & 42.81 & $-$23.14 & 580 & 7.42 & -- & -- \\
30 & 0.050 & 72 & 689 & 176.79 & 58.94 & 38.01 & $-$20.48 & 612 & 6.92 & 191.77 & 78.64 \\
31 & 0.050 & 73 & 692 & 184.70 & 61.49 & 39.23 & $-$20.99 & 579 & 7.49 & -- & -- \\
32 & 0.050 & 74 & 694 & 169.47 & 57.73 & 36.29 & $-$19.57 & 586 & 6.64 & -- & -- \\
33 & 0.050 & 75 & 698 & 167.19 & 58.83 & 38.33 & $-$20.83 & 588 & 6.39 & -- & -- \\
34 & 0.050 & 76 & 705 & 181.22 & 62.31 & 40.23 & $-$21.77 & 615 & 7.25 & -- & -- \\
35 & 0.050 & 77 & 707 & 170.01 & 55.64 & 35.87 & $-$19.26 & 635 & 6.84 & 170.79 & 68.51 \\
36 & 0.050 & 78 & 709 & 174.51 & 57.43 & 36.74 & $-$19.63 & 650 & 7.19 & -- & -- \\
37 & 0.050 & 79 & 710 & 168.58 & 55.79 & 36.42 & $-$19.64 & 600 & 6.68 & -- & -- \\
38 & 0.050 & 80 & 712 & 178.77 & 57.51 & 35.94 & $-$19.39 & 645 & 7.03 & -- & -- \\
39 & 0.050 & 81 & 715 & 172.35 & 52.21 & 32.96 & $-$17.59 & 677 & 6.93 & -- & -- \\
40 & 0.050 & 82 & 715 & 184.34 & 57.17 & 36.79 & $-$19.78 & 638 & 7.34 & 178.29 & 57.24 \\
\bottomrule
\end{tabular}
\caption[Escenario corredor\_alta: episodios de entrenamiento]{Escenario corredor\_alta: parámetros de entrada y métricas de salida por episodio. 40 episodios de entrenamiento con semilla base 42; la fila fijo es el tiempo fijo con la misma semilla base. La política congelada se evalúa cada 5 episodios con una semilla fija ajena al entrenamiento y a la evaluación final.}
\label{tab:corredor_alta-episodios}
\end{table}



Los entrenamientos del módulo U4 se registran en un archivo por
corrida, nombrado según el escenario, la variante y la semilla base:
\begin{center}
\texttt{resultados\_<escenario>\_ippo\_<variante>\_s<base>.csv}
\end{center}
Las tablas siguientes muestran los cinco primeros episodios y uno de
cada diez de cada entrenamiento (variante y semilla base), con la
recompensa de U2 a reloj fijo, la recompensa que vio el agente, la
entropía de la política sobre las decisiones libres, la divergencia
entre la política vieja y la nueva en cada actualización y la
evaluación intermedia con la semilla 999.

% Generado por tablas_tex.py a partir de los CSV del proyecto. No editar a mano.
\begin{table}[H]
\setstretch{1}
\centering
\scriptsize
\setlength{\tabcolsep}{3pt}
\setlength{\tabcolsep}{2pt}
\begin{tabular}{rrrrrrrrrr}
\toprule
Ep. & lr & Retraso & Espera & $r$/dec. & $r$ agente & Cambios & Entropía & KL & \makecell{Eval.\\retraso} \\
 & ($\times 10^{-4}$) & (s/veh) & (s/veh) & (U2) & (U4) & de fase &  &  & (s/veh) \\
\midrule
1 & 10.0 & 89.96 & 33.23 & $-$17.50 & $-$52.16 & 222 & 0.678 & 0.0098 & -- \\
2 & 9.8 & 94.52 & 41.72 & $-$21.96 & $-$53.31 & 198 & 0.689 & 0.0039 & -- \\
3 & 9.7 & 79.21 & 29.84 & $-$15.96 & $-$45.73 & 205 & 0.665 & 0.0109 & -- \\
4 & 9.5 & 29.94 & 13.04 & $-$7.62 & $-$17.31 & 177 & 0.639 & 0.0084 & -- \\
5 & 9.3 & 36.54 & 17.12 & $-$9.86 & $-$21.11 & 166 & 0.621 & 0.0077 & -- \\
10 & 8.5 & 16.71 & 5.43 & $-$3.50 & $-$8.95 & 143 & 0.498 & 0.0180 & 17.75 \\
20 & 6.8 & 15.53 & 4.78 & $-$3.11 & $-$8.29 & 145 & 0.337 & 0.0054 & 15.18 \\
30 & 5.2 & 14.98 & 4.19 & $-$2.76 & $-$8.01 & 162 & 0.299 & 0.0040 & 15.40 \\
40 & 3.5 & 15.03 & 3.98 & $-$2.62 & $-$7.96 & 168 & 0.277 & 0.0025 & 14.91 \\
50 & 1.8 & 15.21 & 4.26 & $-$2.78 & $-$8.08 & 163 & 0.240 & 0.0015 & 14.88 \\
60 & 0.2 & 15.42 & 4.36 & $-$2.86 & $-$8.12 & 153 & 0.229 & 0.0001 & 14.93 \\
\bottomrule
\end{tabular}
\caption[Escenario cruce: entrenamiento IPPO local, semilla base 42]{Escenario cruce: entrenamiento de IPPO local con semilla base 42, 60 episodios (se muestran los cinco primeros y uno de cada diez; la semilla de SUMO de cada episodio es la base más su número). $r$/dec.\ es la recompensa de U2 a reloj fijo; $r$ agente, la que vio el agente por decisión; la entropía se promedia sobre las decisiones libres.}
\label{tab:cruce-ippo-local-42}
\end{table}


% Generado por tablas_tex.py a partir de los CSV del proyecto. No editar a mano.
\begin{table}[H]
\setstretch{1}
\centering
\scriptsize
\setlength{\tabcolsep}{3pt}
\setlength{\tabcolsep}{2pt}
\begin{tabular}{rrrrrrrrrr}
\toprule
Ep. & lr & Retraso & Espera & $r$/dec. & $r$ agente & Cambios & Entropía & KL & \makecell{Eval.\\retraso} \\
 & ($\times 10^{-4}$) & (s/veh) & (s/veh) & (U2) & (U4) & de fase &  &  & (s/veh) \\
\midrule
1 & 10.0 & 30.28 & 11.95 & $-$2.33 & $-$5.42 & 642 & 0.681 & 0.0079 & -- \\
2 & 9.9 & 29.06 & 10.77 & $-$2.11 & $-$5.22 & 687 & 0.652 & 0.0098 & -- \\
3 & 9.9 & 27.25 & 9.65 & $-$1.88 & $-$4.78 & 670 & 0.602 & 0.0118 & -- \\
4 & 9.8 & 26.53 & 9.30 & $-$1.81 & $-$4.58 & 635 & 0.595 & 0.0067 & -- \\
5 & 9.8 & 24.77 & 8.12 & $-$1.59 & $-$4.28 & 671 & 0.568 & 0.0056 & -- \\
10 & 9.6 & 25.13 & 8.63 & $-$1.67 & $-$4.29 & 615 & 0.430 & 0.0033 & 25.32 \\
20 & 9.1 & 24.70 & 8.52 & $-$1.67 & $-$4.15 & 582 & 0.357 & 0.0027 & 25.79 \\
30 & 8.5 & 24.76 & 8.81 & $-$1.61 & $-$4.06 & 593 & 0.267 & 0.0029 & 24.30 \\
40 & 8.1 & 23.88 & 8.13 & $-$1.62 & $-$4.06 & 592 & 0.261 & 0.0019 & 24.59 \\
50 & 7.6 & 23.58 & 7.85 & $-$1.56 & $-$3.98 & 593 & 0.252 & 0.0029 & 24.42 \\
60 & 7.0 & 24.92 & 9.08 & $-$1.75 & $-$4.21 & 547 & 0.219 & 0.0034 & 23.36 \\
70 & 6.5 & 24.01 & 8.37 & $-$1.65 & $-$4.05 & 563 & 0.198 & 0.0038 & 24.08 \\
80 & 6.0 & 24.50 & 8.69 & $-$1.68 & $-$4.12 & 558 & 0.193 & 0.0032 & 23.55 \\
90 & 5.6 & 24.82 & 9.01 & $-$1.76 & $-$4.15 & 554 & 0.217 & 0.0029 & 24.31 \\
100 & 5.0 & 24.53 & 8.75 & $-$1.72 & $-$4.11 & 557 & 0.247 & 0.0026 & 24.42 \\
110 & 4.5 & 24.82 & 9.04 & $-$1.74 & $-$4.18 & 570 & 0.262 & 0.0029 & 24.22 \\
120 & 4.0 & 23.17 & 7.72 & $-$1.53 & $-$3.95 & 610 & 0.234 & 0.0030 & 23.47 \\
130 & 3.6 & 23.64 & 8.28 & $-$1.55 & $-$3.91 & 599 & 0.181 & 0.0022 & 24.12 \\
140 & 3.0 & 23.75 & 8.37 & $-$1.65 & $-$4.00 & 563 & 0.169 & 0.0027 & 23.92 \\
150 & 2.6 & 25.52 & 9.67 & $-$1.72 & $-$4.13 & 568 & 0.195 & 0.0019 & 23.65 \\
160 & 2.0 & 23.50 & 8.05 & $-$1.55 & $-$3.98 & 607 & 0.188 & 0.0015 & 23.16 \\
170 & 1.6 & 23.70 & 8.20 & $-$1.56 & $-$3.97 & 605 & 0.158 & 0.0015 & 23.49 \\
180 & 1.1 & 23.74 & 8.13 & $-$1.63 & $-$4.05 & 581 & 0.154 & 0.0011 & 23.70 \\
190 & 0.6 & 24.06 & 8.49 & $-$1.63 & $-$4.06 & 571 & 0.149 & 0.0009 & 23.42 \\
200 & 0.1 & 23.55 & 8.30 & $-$1.63 & $-$3.98 & 557 & 0.138 & 0.0001 & 23.54 \\
\bottomrule
\end{tabular}
\caption[Escenario red: entrenamiento IPPO local, semilla base 42]{Escenario red: entrenamiento de IPPO local con semilla base 42, 200 episodios (se muestran los cinco primeros y uno de cada diez; la semilla de SUMO de cada episodio es la base más su número). $r$/dec.\ es la recompensa de U2 a reloj fijo; $r$ agente, la que vio el agente por decisión; la entropía se promedia sobre las decisiones libres.}
\label{tab:red-ippo-local-42}
\end{table}

\begin{table}[H]
\setstretch{1}
\centering
\scriptsize
\setlength{\tabcolsep}{3pt}
\setlength{\tabcolsep}{2pt}
\begin{tabular}{rrrrrrrrrr}
\toprule
Ep. & lr & Retraso & Espera & $r$/dec. & $r$ agente & Cambios & Entropía & KL & \makecell{Eval.\\retraso} \\
 & ($\times 10^{-4}$) & (s/veh) & (s/veh) & (U2) & (U4) & de fase &  &  & (s/veh) \\
\midrule
1 & 10.0 & 30.73 & 12.26 & $-$2.40 & $-$5.51 & 640 & 0.677 & 0.0111 & -- \\
2 & 9.9 & 29.77 & 11.67 & $-$2.31 & $-$5.32 & 612 & 0.633 & 0.0102 & -- \\
3 & 9.9 & 27.17 & 9.98 & $-$1.91 & $-$4.66 & 623 & 0.580 & 0.0072 & -- \\
4 & 9.8 & 26.24 & 9.41 & $-$1.82 & $-$4.49 & 598 & 0.564 & 0.0044 & -- \\
5 & 9.8 & 25.86 & 9.26 & $-$1.80 & $-$4.39 & 601 & 0.543 & 0.0037 & -- \\
10 & 9.6 & 25.23 & 8.73 & $-$1.67 & $-$4.24 & 609 & 0.449 & 0.0034 & 25.16 \\
20 & 9.1 & 24.79 & 8.28 & $-$1.59 & $-$4.20 & 646 & 0.366 & 0.0029 & 25.05 \\
30 & 8.5 & 24.26 & 8.23 & $-$1.60 & $-$4.09 & 602 & 0.282 & 0.0027 & 23.59 \\
40 & 8.1 & 24.35 & 8.47 & $-$1.61 & $-$4.05 & 589 & 0.192 & 0.0034 & 23.73 \\
50 & 7.6 & 23.97 & 7.72 & $-$1.51 & $-$4.10 & 644 & 0.183 & 0.0020 & 23.22 \\
60 & 7.0 & 23.47 & 7.49 & $-$1.52 & $-$4.00 & 608 & 0.158 & 0.0028 & 23.14 \\
70 & 6.5 & 23.85 & 7.77 & $-$1.52 & $-$3.98 & 613 & 0.214 & 0.0055 & 23.48 \\
80 & 6.0 & 23.24 & 7.62 & $-$1.50 & $-$3.96 & 616 & 0.208 & 0.0030 & 23.80 \\
90 & 5.6 & 23.92 & 8.03 & $-$1.50 & $-$3.94 & 598 & 0.194 & 0.0030 & 22.73 \\
100 & 5.0 & 23.51 & 7.76 & $-$1.50 & $-$3.87 & 591 & 0.160 & 0.0029 & 22.32 \\
110 & 4.5 & 23.05 & 7.46 & $-$1.49 & $-$3.86 & 592 & 0.137 & 0.0019 & 22.37 \\
120 & 4.0 & 22.86 & 7.29 & $-$1.45 & $-$3.83 & 595 & 0.125 & 0.0019 & 22.23 \\
130 & 3.6 & 23.12 & 7.64 & $-$1.50 & $-$3.87 & 597 & 0.114 & 0.0021 & 22.11 \\
140 & 3.0 & 22.84 & 7.71 & $-$1.46 & $-$3.75 & 593 & 0.110 & 0.0020 & 22.05 \\
150 & 2.6 & 23.26 & 8.05 & $-$1.51 & $-$3.80 & 586 & 0.104 & 0.0010 & 22.27 \\
160 & 2.0 & 24.15 & 8.40 & $-$1.50 & $-$3.89 & 594 & 0.117 & 0.0011 & 22.56 \\
170 & 1.6 & 22.19 & 7.19 & $-$1.40 & $-$3.65 & 593 & 0.119 & 0.0007 & 22.23 \\
180 & 1.1 & 22.74 & 7.39 & $-$1.45 & $-$3.80 & 597 & 0.120 & 0.0007 & 22.73 \\
190 & 0.6 & 23.06 & 7.71 & $-$1.49 & $-$3.82 & 577 & 0.100 & 0.0004 & 22.37 \\
200 & 0.1 & 22.62 & 7.36 & $-$1.41 & $-$3.75 & 582 & 0.091 & 0.0000 & 22.26 \\
\bottomrule
\end{tabular}
\caption[Escenario red: entrenamiento IPPO local, semilla base 2042]{Escenario red: entrenamiento de IPPO local con semilla base 2042, 200 episodios (se muestran los cinco primeros y uno de cada diez; la semilla de SUMO de cada episodio es la base más su número). $r$/dec.\ es la recompensa de U2 a reloj fijo; $r$ agente, la que vio el agente por decisión; la entropía se promedia sobre las decisiones libres.}
\label{tab:red-ippo-local-2042}
\end{table}

\begin{table}[H]
\setstretch{1}
\centering
\scriptsize
\setlength{\tabcolsep}{3pt}
\setlength{\tabcolsep}{2pt}
\begin{tabular}{rrrrrrrrrr}
\toprule
Ep. & lr & Retraso & Espera & $r$/dec. & $r$ agente & Cambios & Entropía & KL & \makecell{Eval.\\retraso} \\
 & ($\times 10^{-4}$) & (s/veh) & (s/veh) & (U2) & (U4) & de fase &  &  & (s/veh) \\
\midrule
1 & 10.0 & 30.11 & 11.79 & $-$2.32 & $-$5.33 & 623 & 0.679 & 0.0096 & -- \\
2 & 9.9 & 27.87 & 9.98 & $-$1.97 & $-$4.90 & 639 & 0.630 & 0.0102 & -- \\
3 & 9.9 & 27.01 & 9.98 & $-$1.91 & $-$4.61 & 592 & 0.599 & 0.0077 & -- \\
4 & 9.8 & 26.97 & 9.50 & $-$1.79 & $-$4.61 & 646 & 0.589 & 0.0046 & -- \\
5 & 9.8 & 25.57 & 8.83 & $-$1.71 & $-$4.40 & 651 & 0.581 & 0.0051 & -- \\
10 & 9.6 & 24.65 & 8.12 & $-$1.56 & $-$4.19 & 665 & 0.473 & 0.0037 & 24.48 \\
20 & 9.1 & 24.53 & 7.83 & $-$1.50 & $-$4.21 & 701 & 0.395 & 0.0030 & 24.78 \\
30 & 8.5 & 26.10 & 9.14 & $-$1.65 & $-$4.30 & 626 & 0.355 & 0.0021 & 24.68 \\
40 & 8.1 & 23.82 & 8.03 & $-$1.49 & $-$3.88 & 602 & 0.273 & 0.0022 & 23.18 \\
50 & 7.6 & 23.30 & 7.84 & $-$1.45 & $-$3.78 & 601 & 0.203 & 0.0020 & 23.13 \\
60 & 7.0 & 23.46 & 7.90 & $-$1.53 & $-$3.91 & 571 & 0.152 & 0.0016 & 22.74 \\
70 & 6.5 & 24.07 & 8.32 & $-$1.58 & $-$3.96 & 578 & 0.156 & 0.0021 & 22.76 \\
80 & 6.0 & 23.35 & 7.98 & $-$1.48 & $-$3.81 & 589 & 0.153 & 0.0015 & 22.95 \\
90 & 5.6 & 24.98 & 9.03 & $-$1.54 & $-$3.99 & 589 & 0.135 & 0.0016 & 23.63 \\
100 & 5.0 & 22.95 & 7.40 & $-$1.42 & $-$3.79 & 604 & 0.130 & 0.0018 & 22.52 \\
110 & 4.5 & 22.74 & 7.43 & $-$1.45 & $-$3.76 & 602 & 0.126 & 0.0020 & 22.57 \\
120 & 4.0 & 23.42 & 7.67 & $-$1.48 & $-$3.85 & 590 & 0.101 & 0.0031 & 22.58 \\
130 & 3.6 & 22.74 & 7.42 & $-$1.45 & $-$3.79 & 601 & 0.115 & 0.0008 & 23.18 \\
140 & 3.0 & 22.19 & 7.25 & $-$1.40 & $-$3.68 & 596 & 0.103 & 0.0011 & 22.51 \\
150 & 2.6 & 21.97 & 7.01 & $-$1.31 & $-$3.58 & 600 & 0.086 & 0.0009 & 22.87 \\
160 & 2.0 & 22.65 & 7.41 & $-$1.42 & $-$3.75 & 601 & 0.088 & 0.0009 & 22.59 \\
170 & 1.6 & 23.50 & 8.16 & $-$1.45 & $-$3.76 & 593 & 0.085 & 0.0011 & 22.06 \\
180 & 1.1 & 22.79 & 7.71 & $-$1.44 & $-$3.72 & 588 & 0.093 & 0.0007 & 22.79 \\
190 & 0.6 & 22.70 & 7.27 & $-$1.43 & $-$3.76 & 588 & 0.085 & 0.0005 & 22.11 \\
200 & 0.1 & 22.65 & 7.28 & $-$1.41 & $-$3.72 & 581 & 0.081 & 0.0002 & 21.76 \\
\bottomrule
\end{tabular}
\caption[Escenario red: entrenamiento IPPO local, semilla base 4042]{Escenario red: entrenamiento de IPPO local con semilla base 4042, 200 episodios (se muestran los cinco primeros y uno de cada diez; la semilla de SUMO de cada episodio es la base más su número). $r$/dec.\ es la recompensa de U2 a reloj fijo; $r$ agente, la que vio el agente por decisión; la entropía se promedia sobre las decisiones libres.}
\label{tab:red-ippo-local-4042}
\end{table}

\begin{table}[H]
\setstretch{1}
\centering
\scriptsize
\setlength{\tabcolsep}{3pt}
\setlength{\tabcolsep}{2pt}
\begin{tabular}{rrrrrrrrrr}
\toprule
Ep. & lr & Retraso & Espera & $r$/dec. & $r$ agente & Cambios & Entropía & KL & \makecell{Eval.\\retraso} \\
 & ($\times 10^{-4}$) & (s/veh) & (s/veh) & (U2) & (U4) & de fase &  &  & (s/veh) \\
\midrule
1 & 10.0 & 30.28 & 11.95 & $-$2.33 & $-$5.42 & 642 & 0.682 & 0.0072 & -- \\
2 & 9.9 & 28.42 & 10.66 & $-$2.02 & $-$4.95 & 658 & 0.643 & 0.0105 & -- \\
3 & 9.9 & 26.11 & 8.89 & $-$1.70 & $-$4.54 & 689 & 0.595 & 0.0090 & -- \\
4 & 9.8 & 25.66 & 8.56 & $-$1.65 & $-$4.43 & 696 & 0.548 & 0.0071 & -- \\
5 & 9.8 & 24.68 & 7.94 & $-$1.55 & $-$4.29 & 706 & 0.495 & 0.0074 & -- \\
10 & 9.6 & 23.21 & 7.45 & $-$1.46 & $-$3.96 & 653 & 0.361 & 0.0021 & 21.96 \\
20 & 9.1 & 22.27 & 7.57 & $-$1.44 & $-$3.68 & 601 & 0.109 & 0.0037 & 21.63 \\
30 & 8.5 & 22.39 & 7.74 & $-$1.41 & $-$3.59 & 593 & 0.213 & 0.0017 & 22.06 \\
40 & 8.1 & 22.12 & 7.54 & $-$1.50 & $-$3.67 & 559 & 0.094 & 0.0008 & 21.65 \\
50 & 7.6 & 22.31 & 7.74 & $-$1.51 & $-$3.71 & 576 & 0.111 & 0.0023 & 19.49 \\
60 & 7.0 & 22.30 & 7.76 & $-$1.50 & $-$3.73 & 573 & 0.060 & 0.0022 & 21.75 \\
70 & 6.5 & 21.88 & 7.48 & $-$1.46 & $-$3.61 & 562 & 0.083 & 0.0013 & 21.65 \\
80 & 6.0 & 22.13 & 7.60 & $-$1.48 & $-$3.64 & 570 & 0.117 & 0.0014 & 21.42 \\
90 & 5.6 & 21.96 & 7.62 & $-$1.48 & $-$3.63 & 572 & 0.089 & 0.0017 & 21.68 \\
100 & 5.0 & 21.78 & 7.24 & $-$1.46 & $-$3.66 & 601 & 0.117 & 0.0017 & 21.68 \\
110 & 4.5 & 22.54 & 8.05 & $-$1.48 & $-$3.65 & 554 & 0.076 & 0.0010 & 21.40 \\
120 & 4.0 & 22.46 & 7.84 & $-$1.48 & $-$3.67 & 551 & 0.081 & 0.0008 & 21.43 \\
130 & 3.6 & 22.79 & 8.12 & $-$1.44 & $-$3.61 & 564 & 0.092 & 0.0007 & 21.87 \\
140 & 3.0 & 22.15 & 7.72 & $-$1.47 & $-$3.63 & 569 & 0.070 & 0.0009 & 21.70 \\
150 & 2.6 & 22.28 & 7.68 & $-$1.41 & $-$3.57 & 558 & 0.114 & 0.0020 & 21.45 \\
160 & 2.0 & 22.33 & 7.69 & $-$1.50 & $-$3.73 & 566 & 0.102 & 0.0014 & 21.84 \\
170 & 1.6 & 21.99 & 7.52 & $-$1.46 & $-$3.64 & 554 & 0.120 & 0.0007 & 22.04 \\
180 & 1.1 & 22.02 & 7.35 & $-$1.44 & $-$3.67 & 550 & 0.110 & 0.0012 & 21.77 \\
190 & 0.6 & 22.39 & 7.85 & $-$1.46 & $-$3.66 & 561 & 0.080 & 0.0008 & 21.74 \\
200 & 0.1 & 21.69 & 7.24 & $-$1.43 & $-$3.58 & 555 & 0.082 & 0.0000 & 21.94 \\
\bottomrule
\end{tabular}
\caption[Escenario red: entrenamiento IPPO vecinos, semilla base 42]{Escenario red: entrenamiento de IPPO vecinos con semilla base 42, 200 episodios (se muestran los cinco primeros y uno de cada diez; la semilla de SUMO de cada episodio es la base más su número). $r$/dec.\ es la recompensa de U2 a reloj fijo; $r$ agente, la que vio el agente por decisión; la entropía se promedia sobre las decisiones libres.}
\label{tab:red-ippo-vecinos-42}
\end{table}

\begin{table}[H]
\setstretch{1}
\centering
\scriptsize
\setlength{\tabcolsep}{3pt}
\setlength{\tabcolsep}{2pt}
\begin{tabular}{rrrrrrrrrr}
\toprule
Ep. & lr & Retraso & Espera & $r$/dec. & $r$ agente & Cambios & Entropía & KL & \makecell{Eval.\\retraso} \\
 & ($\times 10^{-4}$) & (s/veh) & (s/veh) & (U2) & (U4) & de fase &  &  & (s/veh) \\
\midrule
1 & 10.0 & 30.73 & 12.26 & $-$2.40 & $-$5.51 & 640 & 0.676 & 0.0112 & -- \\
2 & 9.9 & 28.92 & 11.22 & $-$2.27 & $-$5.17 & 603 & 0.642 & 0.0096 & -- \\
3 & 9.9 & 25.89 & 8.93 & $-$1.74 & $-$4.45 & 646 & 0.578 & 0.0127 & -- \\
4 & 9.8 & 25.38 & 8.76 & $-$1.72 & $-$4.38 & 622 & 0.533 & 0.0089 & -- \\
5 & 9.8 & 24.55 & 8.37 & $-$1.65 & $-$4.24 & 641 & 0.476 & 0.0084 & -- \\
10 & 9.6 & 23.29 & 7.71 & $-$1.50 & $-$3.94 & 632 & 0.352 & 0.0041 & 21.51 \\
20 & 9.1 & 22.13 & 7.48 & $-$1.49 & $-$3.69 & 576 & 0.184 & 0.0026 & 21.93 \\
30 & 8.5 & 22.29 & 7.71 & $-$1.46 & $-$3.61 & 554 & 0.079 & 0.0014 & 21.23 \\
40 & 8.1 & 21.88 & 7.56 & $-$1.44 & $-$3.57 & 560 & 0.059 & 0.0007 & 21.57 \\
50 & 7.6 & 22.07 & 7.46 & $-$1.48 & $-$3.69 & 557 & 0.096 & 0.0027 & 21.62 \\
60 & 7.0 & 22.41 & 7.97 & $-$1.57 & $-$3.75 & 545 & 0.101 & 0.0020 & 21.45 \\
70 & 6.5 & 22.02 & 7.40 & $-$1.46 & $-$3.66 & 571 & 0.122 & 0.0030 & 21.59 \\
80 & 6.0 & 21.93 & 7.47 & $-$1.49 & $-$3.66 & 540 & 0.126 & 0.0014 & 21.27 \\
90 & 5.6 & 22.02 & 7.63 & $-$1.49 & $-$3.66 & 555 & 0.097 & 0.0020 & 21.75 \\
100 & 5.0 & 21.92 & 7.38 & $-$1.45 & $-$3.59 & 552 & 0.104 & 0.0018 & 21.71 \\
110 & 4.5 & 22.08 & 7.67 & $-$1.48 & $-$3.64 & 558 & 0.123 & 0.0037 & 21.55 \\
120 & 4.0 & 21.74 & 7.53 & $-$1.48 & $-$3.60 & 532 & 0.148 & 0.0022 & 22.01 \\
130 & 3.6 & 22.34 & 7.74 & $-$1.52 & $-$3.70 & 538 & 0.118 & 0.0015 & 21.75 \\
140 & 3.0 & 22.41 & 7.94 & $-$1.50 & $-$3.65 & 537 & 0.129 & 0.0018 & 21.84 \\
150 & 2.6 & 22.27 & 8.06 & $-$1.50 & $-$3.60 & 530 & 0.105 & 0.0008 & 21.52 \\
160 & 2.0 & 23.11 & 8.37 & $-$1.55 & $-$3.74 & 520 & 0.109 & 0.0025 & 22.30 \\
170 & 1.6 & 21.90 & 7.77 & $-$1.45 & $-$3.53 & 526 & 0.095 & 0.0022 & 21.83 \\
180 & 1.1 & 22.27 & 7.69 & $-$1.50 & $-$3.68 & 517 & 0.094 & 0.0014 & 21.88 \\
190 & 0.6 & 21.58 & 7.34 & $-$1.45 & $-$3.57 & 513 & 0.084 & 0.0009 & 21.26 \\
200 & 0.1 & 21.82 & 7.41 & $-$1.48 & $-$3.66 & 520 & 0.088 & 0.0003 & 21.52 \\
\bottomrule
\end{tabular}
\caption[Escenario red: entrenamiento IPPO vecinos, semilla base 2042]{Escenario red: entrenamiento de IPPO vecinos con semilla base 2042, 200 episodios (se muestran los cinco primeros y uno de cada diez; la semilla de SUMO de cada episodio es la base más su número). $r$/dec.\ es la recompensa de U2 a reloj fijo; $r$ agente, la que vio el agente por decisión; la entropía se promedia sobre las decisiones libres.}
\label{tab:red-ippo-vecinos-2042}
\end{table}

\begin{table}[H]
\setstretch{1}
\centering
\scriptsize
\setlength{\tabcolsep}{3pt}
\setlength{\tabcolsep}{2pt}
\begin{tabular}{rrrrrrrrrr}
\toprule
Ep. & lr & Retraso & Espera & $r$/dec. & $r$ agente & Cambios & Entropía & KL & \makecell{Eval.\\retraso} \\
 & ($\times 10^{-4}$) & (s/veh) & (s/veh) & (U2) & (U4) & de fase &  &  & (s/veh) \\
\midrule
1 & 10.0 & 30.54 & 12.00 & $-$2.38 & $-$5.45 & 623 & 0.679 & 0.0095 & -- \\
2 & 9.9 & 28.82 & 11.02 & $-$2.21 & $-$5.13 & 625 & 0.645 & 0.0106 & -- \\
3 & 9.9 & 27.00 & 9.65 & $-$1.88 & $-$4.71 & 649 & 0.595 & 0.0115 & -- \\
4 & 9.8 & 25.69 & 9.00 & $-$1.72 & $-$4.46 & 668 & 0.524 & 0.0146 & -- \\
5 & 9.8 & 25.16 & 8.59 & $-$1.59 & $-$4.30 & 689 & 0.482 & 0.0058 & -- \\
10 & 9.6 & 22.77 & 7.53 & $-$1.45 & $-$3.82 & 625 & 0.318 & 0.0036 & 21.42 \\
20 & 9.1 & 22.59 & 7.79 & $-$1.45 & $-$3.67 & 589 & 0.084 & 0.0027 & 21.42 \\
30 & 8.5 & 22.41 & 7.73 & $-$1.46 & $-$3.66 & 578 & 0.126 & 0.0038 & 21.50 \\
40 & 8.1 & 22.57 & 8.05 & $-$1.44 & $-$3.56 & 582 & 0.016 & 0.0011 & 19.49 \\
50 & 7.6 & 21.90 & 7.24 & $-$1.39 & $-$3.58 & 599 & 0.096 & 0.0051 & 21.23 \\
60 & 7.0 & 22.50 & 7.63 & $-$1.48 & $-$3.74 & 572 & 0.062 & 0.0032 & 21.95 \\
70 & 6.5 & 22.01 & 7.52 & $-$1.40 & $-$3.56 & 581 & 0.039 & 0.0003 & 21.45 \\
80 & 6.0 & 22.73 & 8.13 & $-$1.53 & $-$3.74 & 557 & 0.081 & 0.0029 & 21.68 \\
90 & 5.6 & 22.56 & 8.22 & $-$1.52 & $-$3.70 & 562 & 0.058 & 0.0008 & 21.42 \\
100 & 5.0 & 22.13 & 7.68 & $-$1.49 & $-$3.64 & 563 & 0.055 & 0.0009 & 21.92 \\
110 & 4.5 & 21.86 & 7.25 & $-$1.42 & $-$3.62 & 594 & 0.113 & 0.0013 & 21.59 \\
120 & 4.0 & 21.76 & 7.48 & $-$1.48 & $-$3.61 & 569 & 0.042 & 0.0004 & 21.75 \\
130 & 3.6 & 21.52 & 7.17 & $-$1.41 & $-$3.57 & 563 & 0.108 & 0.0029 & 21.60 \\
140 & 3.0 & 22.16 & 7.54 & $-$1.45 & $-$3.64 & 586 & 0.088 & 0.0008 & 21.59 \\
150 & 2.6 & 22.08 & 7.59 & $-$1.43 & $-$3.60 & 570 & 0.083 & 0.0008 & 21.67 \\
160 & 2.0 & 21.81 & 7.49 & $-$1.44 & $-$3.59 & 565 & 0.065 & 0.0003 & 21.65 \\
170 & 1.6 & 22.40 & 7.88 & $-$1.43 & $-$3.59 & 567 & 0.042 & 0.0004 & 21.83 \\
180 & 1.1 & 22.18 & 7.80 & $-$1.50 & $-$3.65 & 566 & 0.022 & 0.0001 & 21.62 \\
190 & 0.6 & 21.91 & 7.51 & $-$1.42 & $-$3.58 & 572 & 0.039 & 0.0005 & 21.62 \\
200 & 0.1 & 22.16 & 7.63 & $-$1.48 & $-$3.68 & 573 & 0.031 & 0.0000 & 21.66 \\
\bottomrule
\end{tabular}
\caption[Escenario red: entrenamiento IPPO vecinos, semilla base 4042]{Escenario red: entrenamiento de IPPO vecinos con semilla base 4042, 200 episodios (se muestran los cinco primeros y uno de cada diez; la semilla de SUMO de cada episodio es la base más su número). $r$/dec.\ es la recompensa de U2 a reloj fijo; $r$ agente, la que vio el agente por decisión; la entropía se promedia sobre las decisiones libres.}
\label{tab:red-ippo-vecinos-4042}
\end{table}

\begin{table}[H]
\setstretch{1}
\centering
\scriptsize
\setlength{\tabcolsep}{3pt}
\setlength{\tabcolsep}{2pt}
\begin{tabular}{rrrrrrrrrr}
\toprule
Ep. & lr & Retraso & Espera & $r$/dec. & $r$ agente & Cambios & Entropía & KL & \makecell{Eval.\\retraso} \\
 & ($\times 10^{-4}$) & (s/veh) & (s/veh) & (U2) & (U4) & de fase &  &  & (s/veh) \\
\midrule
1 & 10.0 & 30.28 & 11.95 & $-$2.33 & $-$5.42 & 642 & 0.682 & 0.0072 & -- \\
2 & 9.9 & 28.42 & 10.66 & $-$2.02 & $-$4.95 & 658 & 0.643 & 0.0105 & -- \\
3 & 9.9 & 26.11 & 8.89 & $-$1.70 & $-$4.54 & 689 & 0.595 & 0.0090 & -- \\
4 & 9.8 & 25.66 & 8.56 & $-$1.65 & $-$4.43 & 696 & 0.548 & 0.0071 & -- \\
5 & 9.8 & 24.68 & 7.94 & $-$1.55 & $-$4.29 & 706 & 0.495 & 0.0074 & -- \\
10 & 9.6 & 23.21 & 7.45 & $-$1.46 & $-$3.96 & 653 & 0.361 & 0.0021 & 21.96 \\
20 & 9.1 & 22.27 & 7.57 & $-$1.44 & $-$3.68 & 601 & 0.109 & 0.0037 & 21.63 \\
30 & 8.5 & 22.39 & 7.74 & $-$1.41 & $-$3.59 & 593 & 0.213 & 0.0017 & 22.06 \\
40 & 8.1 & 22.12 & 7.54 & $-$1.50 & $-$3.67 & 559 & 0.094 & 0.0008 & 21.65 \\
50 & 7.6 & 22.31 & 7.74 & $-$1.51 & $-$3.71 & 576 & 0.111 & 0.0023 & 19.49 \\
60 & 7.0 & 22.30 & 7.76 & $-$1.50 & $-$3.73 & 573 & 0.060 & 0.0022 & 21.75 \\
70 & 6.5 & 21.88 & 7.48 & $-$1.46 & $-$3.61 & 562 & 0.083 & 0.0013 & 21.65 \\
80 & 6.0 & 22.13 & 7.60 & $-$1.48 & $-$3.64 & 570 & 0.117 & 0.0014 & 21.42 \\
90 & 5.6 & 21.96 & 7.62 & $-$1.48 & $-$3.63 & 572 & 0.089 & 0.0017 & 21.68 \\
100 & 5.0 & 21.78 & 7.24 & $-$1.46 & $-$3.66 & 601 & 0.117 & 0.0017 & 21.68 \\
110 & 4.5 & 22.54 & 8.05 & $-$1.48 & $-$3.65 & 554 & 0.076 & 0.0010 & 21.40 \\
120 & 4.0 & 22.46 & 7.84 & $-$1.48 & $-$3.67 & 551 & 0.081 & 0.0008 & 21.43 \\
130 & 3.6 & 22.79 & 8.12 & $-$1.44 & $-$3.61 & 564 & 0.092 & 0.0007 & 21.87 \\
140 & 3.0 & 22.15 & 7.72 & $-$1.47 & $-$3.63 & 569 & 0.070 & 0.0009 & 21.70 \\
150 & 2.6 & 22.28 & 7.68 & $-$1.41 & $-$3.57 & 558 & 0.114 & 0.0020 & 21.45 \\
160 & 2.0 & 22.33 & 7.69 & $-$1.50 & $-$3.73 & 566 & 0.102 & 0.0014 & 21.84 \\
170 & 1.6 & 21.99 & 7.52 & $-$1.46 & $-$3.64 & 554 & 0.120 & 0.0007 & 22.04 \\
180 & 1.1 & 22.02 & 7.35 & $-$1.44 & $-$3.67 & 550 & 0.110 & 0.0012 & 21.77 \\
190 & 0.6 & 22.39 & 7.85 & $-$1.46 & $-$3.66 & 561 & 0.080 & 0.0008 & 21.74 \\
200 & 0.1 & 21.69 & 7.24 & $-$1.43 & $-$3.58 & 555 & 0.082 & 0.0000 & 21.94 \\
\bottomrule
\end{tabular}
\caption[Escenario red: entrenamiento IPPO spill, semilla base 42]{Escenario red: entrenamiento de IPPO spill con semilla base 42, 200 episodios (se muestran los cinco primeros y uno de cada diez; la semilla de SUMO de cada episodio es la base más su número). $r$/dec.\ es la recompensa de U2 a reloj fijo; $r$ agente, la que vio el agente por decisión; la entropía se promedia sobre las decisiones libres.}
\label{tab:red-ippo-spill-42}
\end{table}


% Generado por tablas_tex.py a partir de los CSV del proyecto. No editar a mano.
\begin{table}[H]
\setstretch{1}
\centering
\scriptsize
\setlength{\tabcolsep}{3pt}
\setlength{\tabcolsep}{2pt}
\begin{tabular}{rrrrrrrrrr}
\toprule
Ep. & lr & Retraso & Espera & $r$/dec. & $r$ agente & Cambios & Entropía & KL & \makecell{Eval.\\retraso} \\
 & ($\times 10^{-4}$) & (s/veh) & (s/veh) & (U2) & (U4) & de fase &  &  & (s/veh) \\
\midrule
1 & 10.0 & 110.42 & 46.09 & $-$11.24 & $-$30.11 & 462 & 0.683 & 0.0071 & -- \\
2 & 9.9 & 69.27 & 29.40 & $-$7.42 & $-$19.32 & 436 & 0.659 & 0.0121 & -- \\
3 & 9.9 & 49.83 & 20.65 & $-$5.17 & $-$13.75 & 440 & 0.610 & 0.0110 & -- \\
4 & 9.8 & 28.77 & 11.61 & $-$3.12 & $-$8.03 & 378 & 0.572 & 0.0062 & -- \\
5 & 9.8 & 26.26 & 10.48 & $-$2.86 & $-$7.26 & 358 & 0.529 & 0.0116 & -- \\
10 & 9.6 & 23.46 & 8.45 & $-$2.26 & $-$6.50 & 383 & 0.389 & 0.0039 & 23.94 \\
20 & 9.1 & 22.84 & 8.47 & $-$2.30 & $-$6.28 & 358 & 0.310 & 0.0031 & 22.43 \\
30 & 8.5 & 22.52 & 8.11 & $-$2.16 & $-$6.18 & 384 & 0.253 & 0.0035 & 21.69 \\
40 & 8.1 & 24.26 & 8.54 & $-$2.26 & $-$6.71 & 406 & 0.224 & 0.0051 & 21.34 \\
50 & 7.6 & 22.35 & 7.38 & $-$1.97 & $-$6.23 & 427 & 0.220 & 0.0048 & 21.41 \\
60 & 7.0 & 21.96 & 7.21 & $-$1.95 & $-$6.06 & 394 & 0.220 & 0.0022 & 21.67 \\
70 & 6.5 & 22.50 & 7.77 & $-$2.07 & $-$6.17 & 398 & 0.231 & 0.0043 & 21.58 \\
80 & 6.0 & 21.99 & 7.59 & $-$2.04 & $-$6.09 & 396 & 0.198 & 0.0039 & 21.92 \\
90 & 5.6 & 21.64 & 7.50 & $-$2.01 & $-$5.95 & 378 & 0.242 & 0.0055 & 21.19 \\
100 & 5.0 & 21.26 & 7.16 & $-$1.93 & $-$5.84 & 392 & 0.195 & 0.0050 & 21.84 \\
110 & 4.5 & 21.70 & 7.68 & $-$2.06 & $-$5.92 & 364 & 0.196 & 0.0025 & 21.12 \\
120 & 4.0 & 21.57 & 7.44 & $-$2.02 & $-$5.87 & 360 & 0.166 & 0.0036 & 21.55 \\
130 & 3.6 & 20.74 & 7.05 & $-$1.93 & $-$5.67 & 365 & 0.139 & 0.0020 & 21.13 \\
140 & 3.0 & 20.95 & 7.27 & $-$1.97 & $-$5.72 & 371 & 0.147 & 0.0019 & 20.54 \\
150 & 2.6 & 21.63 & 7.34 & $-$1.99 & $-$5.91 & 377 & 0.159 & 0.0027 & 20.96 \\
160 & 2.0 & 20.23 & 6.56 & $-$1.78 & $-$5.49 & 371 & 0.134 & 0.0023 & 21.75 \\
170 & 1.6 & 20.80 & 7.05 & $-$1.90 & $-$5.56 & 369 & 0.117 & 0.0014 & 20.21 \\
180 & 1.1 & 21.80 & 7.55 & $-$2.05 & $-$5.99 & 377 & 0.127 & 0.0005 & 20.91 \\
190 & 0.6 & 21.14 & 6.99 & $-$1.88 & $-$5.74 & 379 & 0.121 & 0.0017 & 20.33 \\
200 & 0.1 & 21.88 & 7.83 & $-$2.12 & $-$5.99 & 365 & 0.122 & 0.0003 & 20.50 \\
\bottomrule
\end{tabular}
\caption[Escenario corredor: entrenamiento IPPO local, semilla base 42]{Escenario corredor: entrenamiento de IPPO local con semilla base 42, 200 episodios (se muestran los cinco primeros y uno de cada diez; la semilla de SUMO de cada episodio es la base más su número). $r$/dec.\ es la recompensa de U2 a reloj fijo; $r$ agente, la que vio el agente por decisión; la entropía se promedia sobre las decisiones libres.}
\label{tab:corredor-ippo-local-42}
\end{table}

\begin{table}[H]
\setstretch{1}
\centering
\scriptsize
\setlength{\tabcolsep}{3pt}
\setlength{\tabcolsep}{2pt}
\begin{tabular}{rrrrrrrrrr}
\toprule
Ep. & lr & Retraso & Espera & $r$/dec. & $r$ agente & Cambios & Entropía & KL & \makecell{Eval.\\retraso} \\
 & ($\times 10^{-4}$) & (s/veh) & (s/veh) & (U2) & (U4) & de fase &  &  & (s/veh) \\
\midrule
1 & 10.0 & 75.56 & 31.93 & $-$7.68 & $-$20.23 & 468 & 0.688 & 0.0036 & -- \\
2 & 9.9 & 69.72 & 31.31 & $-$7.77 & $-$18.94 & 432 & 0.663 & 0.0099 & -- \\
3 & 9.9 & 53.25 & 21.35 & $-$5.51 & $-$14.83 & 391 & 0.602 & 0.0112 & -- \\
4 & 9.8 & 32.45 & 13.50 & $-$3.61 & $-$9.05 & 368 & 0.541 & 0.0076 & -- \\
5 & 9.8 & 28.93 & 12.48 & $-$3.40 & $-$8.01 & 341 & 0.516 & 0.0082 & -- \\
10 & 9.6 & 23.18 & 8.52 & $-$2.33 & $-$6.40 & 378 & 0.411 & 0.0053 & 21.88 \\
20 & 9.1 & 22.55 & 8.02 & $-$2.15 & $-$6.21 & 377 & 0.262 & 0.0033 & 22.86 \\
30 & 8.5 & 22.68 & 8.05 & $-$2.18 & $-$6.20 & 360 & 0.233 & 0.0031 & 22.28 \\
40 & 8.1 & 22.46 & 8.09 & $-$2.18 & $-$6.18 & 373 & 0.172 & 0.0026 & 22.67 \\
50 & 7.6 & 22.03 & 7.58 & $-$2.06 & $-$6.07 & 371 & 0.162 & 0.0018 & 22.28 \\
60 & 7.0 & 22.20 & 7.45 & $-$2.02 & $-$6.06 & 381 & 0.207 & 0.0043 & 21.71 \\
70 & 6.5 & 21.86 & 7.24 & $-$1.94 & $-$6.00 & 398 & 0.190 & 0.0026 & 21.85 \\
80 & 6.0 & 21.93 & 7.06 & $-$1.91 & $-$6.10 & 406 & 0.230 & 0.0027 & 22.25 \\
90 & 5.6 & 21.77 & 7.19 & $-$1.93 & $-$6.02 & 401 & 0.190 & 0.0042 & 21.91 \\
100 & 5.0 & 21.42 & 7.26 & $-$1.97 & $-$5.87 & 378 & 0.189 & 0.0028 & 22.44 \\
110 & 4.5 & 21.22 & 7.08 & $-$1.90 & $-$5.83 & 392 & 0.239 & 0.0035 & 22.68 \\
120 & 4.0 & 21.31 & 7.02 & $-$1.89 & $-$5.87 & 396 & 0.203 & 0.0029 & 22.10 \\
130 & 3.6 & 21.75 & 7.38 & $-$1.97 & $-$5.98 & 398 & 0.196 & 0.0029 & 22.10 \\
140 & 3.0 & 22.84 & 7.83 & $-$2.09 & $-$6.30 & 405 & 0.209 & 0.0063 & 22.05 \\
150 & 2.6 & 21.98 & 7.48 & $-$2.00 & $-$6.04 & 393 & 0.201 & 0.0028 & 21.68 \\
160 & 2.0 & 21.53 & 7.36 & $-$2.01 & $-$5.92 & 377 & 0.192 & 0.0023 & 22.68 \\
170 & 1.6 & 21.46 & 7.21 & $-$1.93 & $-$5.93 & 392 & 0.184 & 0.0012 & 21.34 \\
180 & 1.1 & 20.86 & 7.25 & $-$1.94 & $-$5.65 & 371 & 0.138 & 0.0005 & 20.99 \\
190 & 0.6 & 21.27 & 7.14 & $-$1.92 & $-$5.83 & 383 & 0.160 & 0.0008 & 20.55 \\
200 & 0.1 & 20.85 & 6.99 & $-$1.90 & $-$5.71 & 372 & 0.155 & 0.0000 & 20.83 \\
\bottomrule
\end{tabular}
\caption[Escenario corredor: entrenamiento IPPO local, semilla base 2042]{Escenario corredor: entrenamiento de IPPO local con semilla base 2042, 200 episodios (se muestran los cinco primeros y uno de cada diez; la semilla de SUMO de cada episodio es la base más su número). $r$/dec.\ es la recompensa de U2 a reloj fijo; $r$ agente, la que vio el agente por decisión; la entropía se promedia sobre las decisiones libres.}
\label{tab:corredor-ippo-local-2042}
\end{table}

\begin{table}[H]
\setstretch{1}
\centering
\scriptsize
\setlength{\tabcolsep}{3pt}
\setlength{\tabcolsep}{2pt}
\begin{tabular}{rrrrrrrrrr}
\toprule
Ep. & lr & Retraso & Espera & $r$/dec. & $r$ agente & Cambios & Entropía & KL & \makecell{Eval.\\retraso} \\
 & ($\times 10^{-4}$) & (s/veh) & (s/veh) & (U2) & (U4) & de fase &  &  & (s/veh) \\
\midrule
1 & 10.0 & 70.85 & 29.29 & $-$7.42 & $-$19.81 & 431 & 0.683 & 0.0065 & -- \\
2 & 9.9 & 61.59 & 27.91 & $-$7.23 & $-$17.35 & 397 & 0.649 & 0.0130 & -- \\
3 & 9.9 & 36.49 & 15.06 & $-$3.96 & $-$10.02 & 376 & 0.590 & 0.0078 & -- \\
4 & 9.8 & 28.66 & 12.31 & $-$3.35 & $-$7.97 & 349 & 0.569 & 0.0045 & -- \\
5 & 9.8 & 25.18 & 9.85 & $-$2.68 & $-$7.00 & 360 & 0.523 & 0.0071 & -- \\
10 & 9.6 & 24.61 & 9.02 & $-$2.43 & $-$6.78 & 379 & 0.395 & 0.0029 & 23.14 \\
20 & 9.1 & 22.62 & 7.93 & $-$2.16 & $-$6.20 & 377 & 0.318 & 0.0043 & 21.57 \\
30 & 8.5 & 21.97 & 7.60 & $-$2.05 & $-$6.06 & 384 & 0.266 & 0.0014 & 21.45 \\
40 & 8.1 & 23.00 & 7.78 & $-$2.06 & $-$6.34 & 416 & 0.255 & 0.0044 & 21.34 \\
50 & 7.6 & 22.61 & 7.63 & $-$2.01 & $-$6.29 & 429 & 0.211 & 0.0055 & 21.79 \\
60 & 7.0 & 21.88 & 7.13 & $-$1.92 & $-$6.10 & 417 & 0.165 & 0.0063 & 21.51 \\
70 & 6.5 & 21.29 & 6.97 & $-$1.85 & $-$5.84 & 408 & 0.192 & 0.0049 & 21.75 \\
80 & 6.0 & 21.84 & 7.45 & $-$1.96 & $-$6.00 & 392 & 0.241 & 0.0047 & 21.23 \\
90 & 5.6 & 21.95 & 7.32 & $-$1.95 & $-$6.06 & 410 & 0.209 & 0.0045 & 21.84 \\
100 & 5.0 & 21.58 & 7.25 & $-$1.95 & $-$5.98 & 403 & 0.204 & 0.0035 & 21.75 \\
110 & 4.5 & 21.83 & 7.44 & $-$2.01 & $-$6.01 & 390 & 0.204 & 0.0020 & 21.61 \\
120 & 4.0 & 22.31 & 7.64 & $-$2.06 & $-$6.12 & 378 & 0.187 & 0.0019 & 21.65 \\
130 & 3.6 & 20.99 & 7.01 & $-$1.89 & $-$5.68 & 381 & 0.154 & 0.0016 & 20.44 \\
140 & 3.0 & 20.63 & 6.99 & $-$1.90 & $-$5.61 & 375 & 0.144 & 0.0026 & 21.31 \\
150 & 2.6 & 22.65 & 8.19 & $-$2.18 & $-$6.18 & 378 & 0.196 & 0.0030 & 20.58 \\
160 & 2.0 & 21.14 & 7.36 & $-$2.00 & $-$5.82 & 367 & 0.154 & 0.0013 & 21.40 \\
170 & 1.6 & 21.26 & 7.10 & $-$1.91 & $-$5.81 & 387 & 0.157 & 0.0009 & 21.51 \\
180 & 1.1 & 21.57 & 7.48 & $-$2.02 & $-$5.94 & 383 & 0.150 & 0.0010 & 21.37 \\
190 & 0.6 & 21.15 & 7.27 & $-$1.96 & $-$5.74 & 369 & 0.118 & 0.0005 & 21.14 \\
200 & 0.1 & 21.14 & 7.21 & $-$1.95 & $-$5.79 & 379 & 0.120 & 0.0001 & 21.71 \\
\bottomrule
\end{tabular}
\caption[Escenario corredor: entrenamiento IPPO local, semilla base 4042]{Escenario corredor: entrenamiento de IPPO local con semilla base 4042, 200 episodios (se muestran los cinco primeros y uno de cada diez; la semilla de SUMO de cada episodio es la base más su número). $r$/dec.\ es la recompensa de U2 a reloj fijo; $r$ agente, la que vio el agente por decisión; la entropía se promedia sobre las decisiones libres.}
\label{tab:corredor-ippo-local-4042}
\end{table}

\begin{table}[H]
\setstretch{1}
\centering
\scriptsize
\setlength{\tabcolsep}{3pt}
\setlength{\tabcolsep}{2pt}
\begin{tabular}{rrrrrrrrrr}
\toprule
Ep. & lr & Retraso & Espera & $r$/dec. & $r$ agente & Cambios & Entropía & KL & \makecell{Eval.\\retraso} \\
 & ($\times 10^{-4}$) & (s/veh) & (s/veh) & (U2) & (U4) & de fase &  &  & (s/veh) \\
\midrule
1 & 10.0 & 110.42 & 46.09 & $-$11.24 & $-$30.11 & 462 & 0.682 & 0.0075 & -- \\
2 & 9.9 & 68.76 & 28.84 & $-$7.31 & $-$19.24 & 433 & 0.657 & 0.0116 & -- \\
3 & 9.9 & 55.84 & 22.51 & $-$5.63 & $-$15.21 & 432 & 0.627 & 0.0083 & -- \\
4 & 9.8 & 48.97 & 19.34 & $-$4.99 & $-$13.58 & 402 & 0.587 & 0.0091 & -- \\
5 & 9.8 & 26.50 & 10.10 & $-$2.73 & $-$7.43 & 403 & 0.532 & 0.0105 & -- \\
10 & 9.6 & 23.30 & 8.39 & $-$2.28 & $-$6.42 & 370 & 0.383 & 0.0062 & 21.01 \\
20 & 9.1 & 21.59 & 7.57 & $-$2.05 & $-$5.94 & 386 & 0.184 & 0.0064 & 21.14 \\
30 & 8.5 & 21.19 & 7.23 & $-$1.95 & $-$5.82 & 389 & 0.218 & 0.0028 & 21.13 \\
40 & 8.1 & 20.36 & 6.70 & $-$1.82 & $-$5.59 & 384 & 0.132 & 0.0027 & 20.37 \\
50 & 7.6 & 20.09 & 6.69 & $-$1.81 & $-$5.45 & 376 & 0.124 & 0.0029 & 19.98 \\
60 & 7.0 & 21.66 & 7.58 & $-$2.06 & $-$5.89 & 360 & 0.153 & 0.0038 & 19.94 \\
70 & 6.5 & 20.29 & 6.81 & $-$1.83 & $-$5.56 & 386 & 0.118 & 0.0028 & 20.37 \\
80 & 6.0 & 21.07 & 7.12 & $-$1.91 & $-$5.74 & 382 & 0.175 & 0.0058 & 20.43 \\
90 & 5.6 & 20.56 & 7.01 & $-$1.89 & $-$5.65 & 388 & 0.186 & 0.0029 & 19.56 \\
100 & 5.0 & 20.49 & 6.83 & $-$1.82 & $-$5.60 & 401 & 0.121 & 0.0051 & 20.20 \\
110 & 4.5 & 21.13 & 7.41 & $-$2.01 & $-$5.79 & 378 & 0.089 & 0.0042 & 19.44 \\
120 & 4.0 & 20.19 & 6.58 & $-$1.80 & $-$5.53 & 378 & 0.153 & 0.0035 & 19.95 \\
130 & 3.6 & 20.44 & 6.70 & $-$1.81 & $-$5.56 & 378 & 0.151 & 0.0028 & 19.92 \\
140 & 3.0 & 19.82 & 6.56 & $-$1.78 & $-$5.40 & 379 & 0.106 & 0.0011 & 20.30 \\
150 & 2.6 & 20.08 & 6.45 & $-$1.76 & $-$5.47 & 383 & 0.116 & 0.0027 & 19.80 \\
160 & 2.0 & 20.23 & 6.59 & $-$1.81 & $-$5.51 & 376 & 0.101 & 0.0015 & 19.97 \\
170 & 1.6 & 19.64 & 6.29 & $-$1.71 & $-$5.34 & 373 & 0.105 & 0.0009 & 19.94 \\
180 & 1.1 & 20.90 & 7.00 & $-$1.88 & $-$5.67 & 380 & 0.112 & 0.0010 & 20.60 \\
190 & 0.6 & 20.01 & 6.55 & $-$1.77 & $-$5.44 & 381 & 0.078 & 0.0008 & 20.09 \\
200 & 0.1 & 19.62 & 6.45 & $-$1.75 & $-$5.34 & 376 & 0.078 & 0.0000 & 20.23 \\
\bottomrule
\end{tabular}
\caption[Escenario corredor: entrenamiento IPPO vecinos, semilla base 42]{Escenario corredor: entrenamiento de IPPO vecinos con semilla base 42, 200 episodios (se muestran los cinco primeros y uno de cada diez; la semilla de SUMO de cada episodio es la base más su número). $r$/dec.\ es la recompensa de U2 a reloj fijo; $r$ agente, la que vio el agente por decisión; la entropía se promedia sobre las decisiones libres.}
\label{tab:corredor-ippo-vecinos-42}
\end{table}

\begin{table}[H]
\setstretch{1}
\centering
\scriptsize
\setlength{\tabcolsep}{3pt}
\setlength{\tabcolsep}{2pt}
\begin{tabular}{rrrrrrrrrr}
\toprule
Ep. & lr & Retraso & Espera & $r$/dec. & $r$ agente & Cambios & Entropía & KL & \makecell{Eval.\\retraso} \\
 & ($\times 10^{-4}$) & (s/veh) & (s/veh) & (U2) & (U4) & de fase &  &  & (s/veh) \\
\midrule
1 & 10.0 & 78.83 & 33.69 & $-$7.98 & $-$20.86 & 475 & 0.685 & 0.0051 & -- \\
2 & 9.9 & 55.74 & 26.11 & $-$6.62 & $-$15.34 & 422 & 0.659 & 0.0119 & -- \\
3 & 9.9 & 42.66 & 17.26 & $-$4.55 & $-$12.05 & 394 & 0.606 & 0.0110 & -- \\
4 & 9.8 & 34.42 & 13.59 & $-$3.60 & $-$9.61 & 376 & 0.561 & 0.0066 & -- \\
5 & 9.8 & 26.32 & 10.54 & $-$2.88 & $-$7.35 & 361 & 0.521 & 0.0093 & -- \\
10 & 9.6 & 22.60 & 8.22 & $-$2.24 & $-$6.20 & 365 & 0.413 & 0.0060 & 21.41 \\
20 & 9.1 & 20.40 & 6.97 & $-$1.87 & $-$5.56 & 375 & 0.224 & 0.0055 & 20.26 \\
30 & 8.5 & 21.09 & 7.36 & $-$2.00 & $-$5.76 & 372 & 0.170 & 0.0035 & 19.61 \\
40 & 8.1 & 20.61 & 6.99 & $-$1.88 & $-$5.67 & 392 & 0.142 & 0.0027 & 20.70 \\
50 & 7.6 & 21.27 & 7.53 & $-$2.03 & $-$5.86 & 373 & 0.148 & 0.0031 & 20.80 \\
60 & 7.0 & 21.30 & 7.39 & $-$2.02 & $-$5.85 & 378 & 0.149 & 0.0054 & 20.01 \\
70 & 6.5 & 20.89 & 7.03 & $-$1.90 & $-$5.69 & 368 & 0.118 & 0.0038 & 20.22 \\
80 & 6.0 & 21.12 & 7.06 & $-$1.92 & $-$5.81 & 374 & 0.142 & 0.0036 & 20.83 \\
90 & 5.6 & 20.39 & 6.70 & $-$1.82 & $-$5.59 & 382 & 0.135 & 0.0090 & 20.04 \\
100 & 5.0 & 20.41 & 6.78 & $-$1.83 & $-$5.56 & 382 & 0.113 & 0.0028 & 19.90 \\
110 & 4.5 & 20.46 & 6.86 & $-$1.86 & $-$5.60 & 384 & 0.121 & 0.0023 & 19.35 \\
120 & 4.0 & 19.72 & 6.50 & $-$1.75 & $-$5.35 & 378 & 0.110 & 0.0027 & 19.98 \\
130 & 3.6 & 20.81 & 7.14 & $-$1.93 & $-$5.67 & 378 & 0.100 & 0.0029 & 20.24 \\
140 & 3.0 & 19.77 & 6.31 & $-$1.71 & $-$5.36 & 387 & 0.088 & 0.0029 & 19.70 \\
150 & 2.6 & 20.98 & 7.07 & $-$1.91 & $-$5.69 & 368 & 0.130 & 0.0015 & 20.57 \\
160 & 2.0 & 19.98 & 6.45 & $-$1.75 & $-$5.46 & 376 & 0.128 & 0.0007 & 20.29 \\
170 & 1.6 & 19.35 & 6.17 & $-$1.68 & $-$5.26 & 375 & 0.103 & 0.0020 & 20.99 \\
180 & 1.1 & 20.14 & 6.53 & $-$1.76 & $-$5.46 & 384 & 0.092 & 0.0004 & 20.10 \\
190 & 0.6 & 20.27 & 6.70 & $-$1.82 & $-$5.57 & 386 & 0.086 & 0.0005 & 19.80 \\
200 & 0.1 & 20.43 & 6.69 & $-$1.80 & $-$5.56 & 381 & 0.080 & 0.0000 & 19.79 \\
\bottomrule
\end{tabular}
\caption[Escenario corredor: entrenamiento IPPO vecinos, semilla base 2042]{Escenario corredor: entrenamiento de IPPO vecinos con semilla base 2042, 200 episodios (se muestran los cinco primeros y uno de cada diez; la semilla de SUMO de cada episodio es la base más su número). $r$/dec.\ es la recompensa de U2 a reloj fijo; $r$ agente, la que vio el agente por decisión; la entropía se promedia sobre las decisiones libres.}
\label{tab:corredor-ippo-vecinos-2042}
\end{table}

\begin{table}[H]
\setstretch{1}
\centering
\scriptsize
\setlength{\tabcolsep}{3pt}
\setlength{\tabcolsep}{2pt}
\begin{tabular}{rrrrrrrrrr}
\toprule
Ep. & lr & Retraso & Espera & $r$/dec. & $r$ agente & Cambios & Entropía & KL & \makecell{Eval.\\retraso} \\
 & ($\times 10^{-4}$) & (s/veh) & (s/veh) & (U2) & (U4) & de fase &  &  & (s/veh) \\
\midrule
1 & 10.0 & 82.44 & 34.77 & $-$8.43 & $-$22.07 & 455 & 0.680 & 0.0086 & -- \\
2 & 9.9 & 58.68 & 27.43 & $-$6.97 & $-$16.21 & 402 & 0.660 & 0.0119 & -- \\
3 & 9.9 & 38.39 & 15.27 & $-$4.09 & $-$10.90 & 374 & 0.603 & 0.0094 & -- \\
4 & 9.8 & 34.37 & 14.56 & $-$3.92 & $-$9.55 & 341 & 0.546 & 0.0106 & -- \\
5 & 9.8 & 26.58 & 11.00 & $-$2.99 & $-$7.41 & 356 & 0.479 & 0.0156 & -- \\
10 & 9.6 & 22.60 & 8.27 & $-$2.23 & $-$6.17 & 372 & 0.362 & 0.0061 & 22.51 \\
20 & 9.1 & 20.85 & 7.30 & $-$1.98 & $-$5.67 & 366 & 0.159 & 0.0019 & 20.88 \\
30 & 8.5 & 20.95 & 7.32 & $-$1.98 & $-$5.69 & 362 & 0.142 & 0.0092 & 19.35 \\
40 & 8.1 & 21.61 & 7.51 & $-$2.02 & $-$5.96 & 390 & 0.141 & 0.0032 & 21.36 \\
50 & 7.6 & 20.69 & 6.93 & $-$1.88 & $-$5.69 & 383 & 0.161 & 0.0058 & 20.31 \\
60 & 7.0 & 23.56 & 8.22 & $-$2.22 & $-$6.55 & 396 & 0.168 & 0.0044 & 21.08 \\
70 & 6.5 & 21.54 & 7.40 & $-$1.99 & $-$5.89 & 384 & 0.138 & 0.0038 & 19.85 \\
80 & 6.0 & 20.21 & 6.72 & $-$1.81 & $-$5.46 & 368 & 0.057 & 0.0009 & 19.79 \\
90 & 5.6 & 20.45 & 6.79 & $-$1.83 & $-$5.56 & 377 & 0.097 & 0.0038 & 19.90 \\
100 & 5.0 & 20.48 & 6.82 & $-$1.84 & $-$5.55 & 373 & 0.105 & 0.0020 & 20.40 \\
110 & 4.5 & 20.82 & 7.06 & $-$1.90 & $-$5.64 & 368 & 0.124 & 0.0031 & 20.04 \\
120 & 4.0 & 20.50 & 6.81 & $-$1.83 & $-$5.60 & 387 & 0.117 & 0.0027 & 20.66 \\
130 & 3.6 & 20.18 & 6.61 & $-$1.78 & $-$5.49 & 388 & 0.121 & 0.0042 & 20.26 \\
140 & 3.0 & 20.14 & 6.65 & $-$1.81 & $-$5.51 & 387 & 0.126 & 0.0025 & 20.43 \\
150 & 2.6 & 19.83 & 6.35 & $-$1.72 & $-$5.40 & 386 & 0.121 & 0.0028 & 20.04 \\
160 & 2.0 & 20.48 & 6.58 & $-$1.78 & $-$5.60 & 391 & 0.112 & 0.0022 & 20.65 \\
170 & 1.6 & 19.76 & 6.24 & $-$1.67 & $-$5.33 & 388 & 0.106 & 0.0015 & 20.43 \\
180 & 1.1 & 19.54 & 5.99 & $-$1.64 & $-$5.34 & 388 & 0.122 & 0.0014 & 20.30 \\
190 & 0.6 & 19.33 & 5.99 & $-$1.62 & $-$5.25 & 389 & 0.101 & 0.0017 & 19.91 \\
200 & 0.1 & 20.11 & 6.48 & $-$1.73 & $-$5.43 & 379 & 0.107 & 0.0001 & 20.70 \\
\bottomrule
\end{tabular}
\caption[Escenario corredor: entrenamiento IPPO vecinos, semilla base 4042]{Escenario corredor: entrenamiento de IPPO vecinos con semilla base 4042, 200 episodios (se muestran los cinco primeros y uno de cada diez; la semilla de SUMO de cada episodio es la base más su número). $r$/dec.\ es la recompensa de U2 a reloj fijo; $r$ agente, la que vio el agente por decisión; la entropía se promedia sobre las decisiones libres.}
\label{tab:corredor-ippo-vecinos-4042}
\end{table}

\begin{table}[H]
\setstretch{1}
\centering
\scriptsize
\setlength{\tabcolsep}{3pt}
\setlength{\tabcolsep}{2pt}
\begin{tabular}{rrrrrrrrrr}
\toprule
Ep. & lr & Retraso & Espera & $r$/dec. & $r$ agente & Cambios & Entropía & KL & \makecell{Eval.\\retraso} \\
 & ($\times 10^{-4}$) & (s/veh) & (s/veh) & (U2) & (U4) & de fase &  &  & (s/veh) \\
\midrule
1 & 10.0 & 110.42 & 46.09 & $-$11.24 & $-$30.37 & 462 & 0.682 & 0.0077 & -- \\
2 & 9.9 & 70.58 & 30.09 & $-$7.58 & $-$19.70 & 436 & 0.657 & 0.0113 & -- \\
3 & 9.9 & 54.94 & 22.45 & $-$5.73 & $-$15.33 & 429 & 0.611 & 0.0111 & -- \\
4 & 9.8 & 37.87 & 15.02 & $-$4.00 & $-$10.67 & 379 & 0.553 & 0.0087 & -- \\
5 & 9.8 & 29.57 & 12.39 & $-$3.44 & $-$8.30 & 355 & 0.557 & 0.0132 & -- \\
10 & 9.6 & 22.78 & 8.16 & $-$2.22 & $-$6.27 & 362 & 0.363 & 0.0060 & 21.26 \\
20 & 9.1 & 21.40 & 7.61 & $-$2.05 & $-$5.88 & 384 & 0.137 & 0.0036 & 22.40 \\
30 & 8.5 & 21.11 & 7.33 & $-$1.97 & $-$5.76 & 369 & 0.187 & 0.0029 & 20.61 \\
40 & 8.1 & 21.25 & 7.46 & $-$2.02 & $-$5.85 & 372 & 0.138 & 0.0022 & 20.01 \\
50 & 7.6 & 21.35 & 7.59 & $-$2.04 & $-$5.84 & 372 & 0.138 & 0.0039 & 20.93 \\
60 & 7.0 & 20.98 & 6.99 & $-$1.88 & $-$5.72 & 372 & 0.077 & 0.0022 & 19.69 \\
70 & 6.5 & 20.79 & 7.07 & $-$1.89 & $-$5.66 & 382 & 0.111 & 0.0026 & 19.99 \\
80 & 6.0 & 20.40 & 6.86 & $-$1.87 & $-$5.59 & 378 & 0.102 & 0.0029 & 20.45 \\
90 & 5.6 & 20.01 & 6.82 & $-$1.85 & $-$5.42 & 366 & 0.061 & 0.0017 & 20.69 \\
100 & 5.0 & 21.06 & 7.50 & $-$2.03 & $-$5.75 & 370 & 0.080 & 0.0030 & 21.20 \\
110 & 4.5 & 20.61 & 7.05 & $-$1.91 & $-$5.62 & 372 & 0.101 & 0.0041 & 19.29 \\
120 & 4.0 & 20.24 & 6.56 & $-$1.77 & $-$5.51 & 386 & 0.091 & 0.0007 & 20.85 \\
130 & 3.6 & 20.53 & 6.72 & $-$1.81 & $-$5.61 & 391 & 0.099 & 0.0015 & 20.39 \\
140 & 3.0 & 19.71 & 6.51 & $-$1.76 & $-$5.35 & 378 & 0.087 & 0.0030 & 21.63 \\
150 & 2.6 & 19.81 & 6.46 & $-$1.75 & $-$5.42 & 381 & 0.081 & 0.0019 & 19.79 \\
160 & 2.0 & 21.71 & 7.51 & $-$2.02 & $-$5.98 & 396 & 0.092 & 0.0019 & 20.51 \\
170 & 1.6 & 20.91 & 7.16 & $-$1.94 & $-$5.70 & 379 & 0.086 & 0.0009 & 20.16 \\
180 & 1.1 & 21.57 & 7.55 & $-$2.02 & $-$5.86 & 373 & 0.067 & 0.0012 & 19.86 \\
190 & 0.6 & 19.78 & 6.41 & $-$1.73 & $-$5.38 & 371 & 0.046 & 0.0005 & 21.21 \\
200 & 0.1 & 19.83 & 6.62 & $-$1.80 & $-$5.37 & 368 & 0.048 & 0.0001 & 20.83 \\
\bottomrule
\end{tabular}
\caption[Escenario corredor: entrenamiento IPPO spill, semilla base 42]{Escenario corredor: entrenamiento de IPPO spill con semilla base 42, 200 episodios (se muestran los cinco primeros y uno de cada diez; la semilla de SUMO de cada episodio es la base más su número). $r$/dec.\ es la recompensa de U2 a reloj fijo; $r$ agente, la que vio el agente por decisión; la entropía se promedia sobre las decisiones libres.}
\label{tab:corredor-ippo-spill-42}
\end{table}


% Generado por tablas_tex.py a partir de los CSV del proyecto. No editar a mano.
\begin{table}[H]
\setstretch{1}
\centering
\scriptsize
\setlength{\tabcolsep}{3pt}
\setlength{\tabcolsep}{2pt}
\begin{tabular}{rrrrrrrrrr}
\toprule
Ep. & lr & Retraso & Espera & $r$/dec. & $r$ agente & Cambios & Entropía & KL & \makecell{Eval.\\retraso} \\
 & ($\times 10^{-4}$) & (s/veh) & (s/veh) & (U2) & (U4) & de fase &  &  & (s/veh) \\
\midrule
1 & 10.0 & 207.25 & 80.36 & $-$26.07 & $-$75.52 & 674 & 0.679 & 0.0095 & -- \\
2 & 9.9 & 198.28 & 85.96 & $-$28.40 & $-$71.94 & 611 & 0.658 & 0.0112 & -- \\
3 & 9.9 & 180.77 & 74.11 & $-$24.75 & $-$66.00 & 596 & 0.625 & 0.0058 & -- \\
4 & 9.8 & 178.00 & 78.09 & $-$26.44 & $-$64.72 & 555 & 0.577 & 0.0069 & -- \\
5 & 9.8 & 172.11 & 86.14 & $-$32.59 & $-$66.98 & 433 & 0.518 & 0.0089 & -- \\
10 & 9.6 & 162.71 & 81.61 & $-$31.19 & $-$63.50 & 357 & 0.374 & 0.0044 & 152.79 \\
20 & 9.1 & 139.13 & 74.94 & $-$29.38 & $-$54.90 & 301 & 0.270 & 0.0032 & 143.76 \\
30 & 8.5 & 144.13 & 76.18 & $-$29.16 & $-$55.70 & 307 & 0.252 & 0.0057 & 132.79 \\
40 & 8.1 & 128.75 & 66.34 & $-$25.62 & $-$50.23 & 327 & 0.237 & 0.0061 & 133.01 \\
50 & 7.6 & 122.54 & 63.22 & $-$23.99 & $-$47.23 & 325 & 0.211 & 0.0047 & 124.80 \\
60 & 7.0 & 137.06 & 69.24 & $-$26.21 & $-$52.80 & 346 & 0.231 & 0.0033 & 127.76 \\
70 & 6.5 & 128.84 & 64.72 & $-$25.11 & $-$50.57 & 323 & 0.211 & 0.0029 & 136.94 \\
80 & 6.0 & 124.43 & 61.69 & $-$23.68 & $-$48.57 & 337 & 0.220 & 0.0025 & 126.83 \\
90 & 5.6 & 122.37 & 62.08 & $-$24.13 & $-$48.31 & 323 & 0.196 & 0.0023 & 136.89 \\
100 & 5.0 & 121.75 & 59.45 & $-$22.86 & $-$47.98 & 345 & 0.201 & 0.0032 & 121.43 \\
110 & 4.5 & 129.19 & 65.93 & $-$25.41 & $-$50.39 & 324 & 0.177 & 0.0027 & 119.69 \\
120 & 4.0 & 121.90 & 63.46 & $-$24.73 & $-$48.17 & 309 & 0.152 & 0.0025 & 125.34 \\
130 & 3.6 & 130.00 & 65.59 & $-$25.06 & $-$50.52 & 329 & 0.189 & 0.0023 & 125.94 \\
140 & 3.0 & 122.89 & 60.91 & $-$23.36 & $-$47.65 & 309 & 0.183 & 0.0028 & 118.87 \\
150 & 2.6 & 122.24 & 65.51 & $-$25.68 & $-$48.19 & 299 & 0.164 & 0.0018 & 117.36 \\
160 & 2.0 & 121.15 & 61.43 & $-$24.26 & $-$48.17 & 310 & 0.158 & 0.0011 & 125.26 \\
170 & 1.6 & 119.30 & 60.88 & $-$23.71 & $-$47.11 & 315 & 0.170 & 0.0011 & 122.99 \\
180 & 1.1 & 129.14 & 66.92 & $-$25.74 & $-$49.97 & 313 & 0.156 & 0.0007 & 122.34 \\
190 & 0.6 & 121.71 & 62.54 & $-$24.28 & $-$47.57 & 307 & 0.156 & 0.0010 & 120.98 \\
200 & 0.1 & 129.71 & 67.24 & $-$26.53 & $-$51.33 & 303 & 0.162 & 0.0001 & 125.37 \\
\bottomrule
\end{tabular}
\caption[Escenario corredor\_alta: entrenamiento IPPO local, semilla base 42]{Escenario corredor\_alta: entrenamiento de IPPO local con semilla base 42, 200 episodios (se muestran los cinco primeros y uno de cada diez; la semilla de SUMO de cada episodio es la base más su número). $r$/dec.\ es la recompensa de U2 a reloj fijo; $r$ agente, la que vio el agente por decisión; la entropía se promedia sobre las decisiones libres.}
\label{tab:corredor_alta-ippo-local-42}
\end{table}

\begin{table}[H]
\setstretch{1}
\centering
\scriptsize
\setlength{\tabcolsep}{3pt}
\setlength{\tabcolsep}{2pt}
\begin{tabular}{rrrrrrrrrr}
\toprule
Ep. & lr & Retraso & Espera & $r$/dec. & $r$ agente & Cambios & Entropía & KL & \makecell{Eval.\\retraso} \\
 & ($\times 10^{-4}$) & (s/veh) & (s/veh) & (U2) & (U4) & de fase &  &  & (s/veh) \\
\midrule
1 & 10.0 & 207.26 & 80.12 & $-$25.94 & $-$75.45 & 674 & 0.681 & 0.0080 & -- \\
2 & 9.9 & 192.11 & 84.21 & $-$27.77 & $-$69.73 & 628 & 0.665 & 0.0100 & -- \\
3 & 9.9 & 186.60 & 75.52 & $-$25.25 & $-$68.14 & 599 & 0.633 & 0.0049 & -- \\
4 & 9.8 & 180.84 & 78.20 & $-$26.31 & $-$66.09 & 583 & 0.589 & 0.0080 & -- \\
5 & 9.8 & 205.09 & 95.34 & $-$34.28 & $-$78.29 & 500 & 0.557 & 0.0081 & -- \\
10 & 9.6 & 155.95 & 78.40 & $-$29.93 & $-$60.95 & 369 & 0.380 & 0.0057 & 153.87 \\
20 & 9.1 & 136.53 & 74.61 & $-$28.98 & $-$53.15 & 317 & 0.275 & 0.0060 & 131.56 \\
30 & 8.5 & 129.37 & 70.64 & $-$27.25 & $-$49.86 & 301 & 0.226 & 0.0062 & 123.08 \\
40 & 8.1 & 131.74 & 66.93 & $-$25.91 & $-$51.32 & 308 & 0.245 & 0.0084 & 135.77 \\
50 & 7.6 & 132.93 & 66.12 & $-$25.27 & $-$51.30 & 321 & 0.221 & 0.0048 & 130.36 \\
60 & 7.0 & 140.27 & 76.75 & $-$29.62 & $-$54.08 & 300 & 0.200 & 0.0052 & 142.05 \\
70 & 6.5 & 133.41 & 69.96 & $-$27.04 & $-$52.02 & 315 & 0.191 & 0.0055 & 131.59 \\
80 & 6.0 & 131.02 & 69.26 & $-$26.36 & $-$50.54 & 318 & 0.190 & 0.0049 & 122.51 \\
90 & 5.6 & 128.98 & 69.01 & $-$26.09 & $-$49.58 & 326 & 0.180 & 0.0029 & 125.88 \\
100 & 5.0 & 130.66 & 65.65 & $-$25.06 & $-$50.42 & 321 & 0.189 & 0.0032 & 124.43 \\
110 & 4.5 & 133.60 & 66.17 & $-$25.47 & $-$52.18 & 325 & 0.161 & 0.0035 & 123.65 \\
120 & 4.0 & 122.32 & 62.50 & $-$24.05 & $-$47.51 & 317 & 0.164 & 0.0028 & 126.68 \\
130 & 3.6 & 124.59 & 66.70 & $-$25.98 & $-$48.66 & 315 & 0.174 & 0.0029 & 133.05 \\
140 & 3.0 & 124.77 & 62.62 & $-$24.60 & $-$49.69 & 324 & 0.198 & 0.0022 & 122.20 \\
150 & 2.6 & 131.59 & 68.38 & $-$26.34 & $-$51.39 & 323 & 0.185 & 0.0025 & 124.14 \\
160 & 2.0 & 134.09 & 69.05 & $-$27.09 & $-$53.03 & 320 & 0.195 & 0.0023 & 123.57 \\
170 & 1.6 & 124.68 & 63.14 & $-$24.06 & $-$47.95 & 319 & 0.174 & 0.0015 & 120.91 \\
180 & 1.1 & 124.21 & 64.79 & $-$25.05 & $-$48.30 & 307 & 0.161 & 0.0014 & 118.24 \\
190 & 0.6 & 125.87 & 66.41 & $-$25.77 & $-$48.87 & 304 & 0.154 & 0.0007 & 119.90 \\
200 & 0.1 & 125.60 & 65.61 & $-$26.22 & $-$50.13 & 315 & 0.153 & 0.0001 & 121.57 \\
\bottomrule
\end{tabular}
\caption[Escenario corredor\_alta: entrenamiento IPPO vecinos, semilla base 42]{Escenario corredor\_alta: entrenamiento de IPPO vecinos con semilla base 42, 200 episodios (se muestran los cinco primeros y uno de cada diez; la semilla de SUMO de cada episodio es la base más su número). $r$/dec.\ es la recompensa de U2 a reloj fijo; $r$ agente, la que vio el agente por decisión; la entropía se promedia sobre las decisiones libres.}
\label{tab:corredor_alta-ippo-vecinos-42}
\end{table}

\begin{table}[H]
\setstretch{1}
\centering
\scriptsize
\setlength{\tabcolsep}{3pt}
\setlength{\tabcolsep}{2pt}
\begin{tabular}{rrrrrrrrrr}
\toprule
Ep. & lr & Retraso & Espera & $r$/dec. & $r$ agente & Cambios & Entropía & KL & \makecell{Eval.\\retraso} \\
 & ($\times 10^{-4}$) & (s/veh) & (s/veh) & (U2) & (U4) & de fase &  &  & (s/veh) \\
\midrule
1 & 10.0 & 204.01 & 79.96 & $-$25.63 & $-$73.75 & 681 & 0.679 & 0.0101 & -- \\
2 & 9.9 & 194.42 & 84.86 & $-$28.32 & $-$70.97 & 600 & 0.655 & 0.0064 & -- \\
3 & 9.9 & 199.18 & 93.98 & $-$31.17 & $-$71.69 & 574 & 0.626 & 0.0101 & -- \\
4 & 9.8 & 170.97 & 78.63 & $-$26.91 & $-$62.46 & 544 & 0.592 & 0.0080 & -- \\
5 & 9.8 & 178.50 & 73.38 & $-$25.07 & $-$65.83 & 553 & 0.565 & 0.0096 & -- \\
10 & 9.6 & 155.23 & 73.41 & $-$27.38 & $-$60.22 & 411 & 0.442 & 0.0061 & 137.58 \\
20 & 9.1 & 134.14 & 71.05 & $-$27.15 & $-$52.03 & 333 & 0.293 & 0.0067 & 143.43 \\
30 & 8.5 & 131.38 & 70.18 & $-$26.93 & $-$50.96 & 332 & 0.264 & 0.0051 & 146.93 \\
40 & 8.1 & 134.57 & 71.91 & $-$27.39 & $-$51.79 & 334 & 0.225 & 0.0051 & 122.71 \\
50 & 7.6 & 120.20 & 61.80 & $-$23.94 & $-$46.74 & 308 & 0.215 & 0.0052 & 121.69 \\
60 & 7.0 & 120.94 & 64.68 & $-$25.03 & $-$47.02 & 299 & 0.215 & 0.0047 & 121.78 \\
70 & 6.5 & 129.67 & 67.47 & $-$26.51 & $-$51.30 & 312 & 0.202 & 0.0045 & 130.77 \\
80 & 6.0 & 129.18 & 67.64 & $-$27.23 & $-$51.30 & 300 & 0.202 & 0.0043 & 123.36 \\
90 & 5.6 & 128.15 & 67.10 & $-$25.97 & $-$49.73 & 303 & 0.186 & 0.0039 & 117.45 \\
100 & 5.0 & 127.56 & 64.30 & $-$25.14 & $-$50.60 & 335 & 0.216 & 0.0040 & 127.18 \\
110 & 4.5 & 119.92 & 62.45 & $-$24.31 & $-$47.22 & 323 & 0.206 & 0.0041 & 122.33 \\
120 & 4.0 & 130.92 & 66.97 & $-$25.93 & $-$51.47 & 330 & 0.201 & 0.0030 & 117.30 \\
130 & 3.6 & 127.44 & 64.14 & $-$24.69 & $-$49.55 & 309 & 0.187 & 0.0024 & 119.54 \\
140 & 3.0 & 115.66 & 61.52 & $-$24.14 & $-$45.67 & 299 & 0.180 & 0.0062 & 117.82 \\
150 & 2.6 & 120.23 & 62.96 & $-$24.37 & $-$46.89 & 302 & 0.150 & 0.0027 & 120.82 \\
160 & 2.0 & 120.34 & 64.06 & $-$25.99 & $-$48.26 & 301 & 0.155 & 0.0011 & 119.04 \\
170 & 1.6 & 121.77 & 63.12 & $-$24.81 & $-$47.86 & 287 & 0.160 & 0.0015 & 123.82 \\
180 & 1.1 & 122.76 & 64.37 & $-$24.82 & $-$47.79 & 303 & 0.161 & 0.0010 & 121.77 \\
190 & 0.6 & 127.79 & 65.84 & $-$25.61 & $-$49.82 & 301 & 0.167 & 0.0011 & 126.17 \\
200 & 0.1 & 118.12 & 60.02 & $-$23.31 & $-$46.11 & 291 & 0.158 & 0.0000 & 119.04 \\
\bottomrule
\end{tabular}
\caption[Escenario corredor\_alta: entrenamiento IPPO vecinos, semilla base 2042]{Escenario corredor\_alta: entrenamiento de IPPO vecinos con semilla base 2042, 200 episodios (se muestran los cinco primeros y uno de cada diez; la semilla de SUMO de cada episodio es la base más su número). $r$/dec.\ es la recompensa de U2 a reloj fijo; $r$ agente, la que vio el agente por decisión; la entropía se promedia sobre las decisiones libres.}
\label{tab:corredor_alta-ippo-vecinos-2042}
\end{table}

\begin{table}[H]
\setstretch{1}
\centering
\scriptsize
\setlength{\tabcolsep}{3pt}
\setlength{\tabcolsep}{2pt}
\begin{tabular}{rrrrrrrrrr}
\toprule
Ep. & lr & Retraso & Espera & $r$/dec. & $r$ agente & Cambios & Entropía & KL & \makecell{Eval.\\retraso} \\
 & ($\times 10^{-4}$) & (s/veh) & (s/veh) & (U2) & (U4) & de fase &  &  & (s/veh) \\
\midrule
1 & 10.0 & 201.61 & 82.22 & $-$26.88 & $-$73.70 & 660 & 0.684 & 0.0057 & -- \\
2 & 9.9 & 199.28 & 80.31 & $-$26.56 & $-$72.92 & 637 & 0.665 & 0.0069 & -- \\
3 & 9.9 & 174.02 & 71.31 & $-$24.12 & $-$64.07 & 586 & 0.615 & 0.0117 & -- \\
4 & 9.8 & 168.85 & 71.79 & $-$25.19 & $-$63.75 & 552 & 0.540 & 0.0157 & -- \\
5 & 9.8 & 199.01 & 92.54 & $-$34.40 & $-$77.44 & 446 & 0.501 & 0.0068 & -- \\
10 & 9.6 & 153.00 & 73.57 & $-$27.61 & $-$58.89 & 361 & 0.355 & 0.0071 & 138.03 \\
20 & 9.1 & 139.60 & 72.37 & $-$27.85 & $-$53.79 & 297 & 0.250 & 0.0062 & 137.62 \\
30 & 8.5 & 130.05 & 68.91 & $-$26.18 & $-$49.64 & 305 & 0.206 & 0.0053 & 131.35 \\
40 & 8.1 & 125.77 & 63.67 & $-$24.44 & $-$49.15 & 335 & 0.218 & 0.0083 & 126.37 \\
50 & 7.6 & 126.56 & 65.68 & $-$24.84 & $-$48.58 & 338 & 0.241 & 0.0057 & 128.74 \\
60 & 7.0 & 123.94 & 63.81 & $-$24.48 & $-$48.29 & 327 & 0.230 & 0.0062 & 123.11 \\
70 & 6.5 & 138.98 & 73.66 & $-$28.41 & $-$54.28 & 315 & 0.218 & 0.0069 & 125.53 \\
80 & 6.0 & 123.77 & 63.79 & $-$24.58 & $-$48.42 & 323 & 0.175 & 0.0044 & 121.70 \\
90 & 5.6 & 124.76 & 62.80 & $-$24.41 & $-$48.96 & 313 & 0.193 & 0.0049 & 127.58 \\
100 & 5.0 & 129.91 & 64.93 & $-$24.69 & $-$50.07 & 319 & 0.189 & 0.0034 & 121.24 \\
110 & 4.5 & 131.47 & 65.43 & $-$24.50 & $-$49.93 & 331 & 0.191 & 0.0037 & 123.39 \\
120 & 4.0 & 124.15 & 62.30 & $-$23.63 & $-$47.69 & 321 & 0.168 & 0.0030 & 124.24 \\
130 & 3.6 & 130.79 & 67.22 & $-$26.11 & $-$51.01 & 314 & 0.165 & 0.0030 & 123.15 \\
140 & 3.0 & 124.41 & 64.16 & $-$24.68 & $-$48.27 & 317 & 0.163 & 0.0033 & 121.70 \\
150 & 2.6 & 123.42 & 61.81 & $-$23.75 & $-$47.91 & 310 & 0.160 & 0.0022 & 123.93 \\
160 & 2.0 & 120.83 & 59.37 & $-$22.66 & $-$46.73 & 316 & 0.154 & 0.0018 & 118.58 \\
170 & 1.6 & 122.21 & 61.43 & $-$23.28 & $-$46.73 & 305 & 0.134 & 0.0014 & 121.38 \\
180 & 1.1 & 126.56 & 67.57 & $-$26.78 & $-$49.85 & 309 & 0.148 & 0.0013 & 120.46 \\
190 & 0.6 & 121.02 & 60.50 & $-$23.17 & $-$46.73 & 311 & 0.129 & 0.0008 & 118.42 \\
200 & 0.1 & 126.77 & 67.92 & $-$27.02 & $-$49.76 & 296 & 0.134 & 0.0001 & 117.50 \\
\bottomrule
\end{tabular}
\caption[Escenario corredor\_alta: entrenamiento IPPO vecinos, semilla base 4042]{Escenario corredor\_alta: entrenamiento de IPPO vecinos con semilla base 4042, 200 episodios (se muestran los cinco primeros y uno de cada diez; la semilla de SUMO de cada episodio es la base más su número). $r$/dec.\ es la recompensa de U2 a reloj fijo; $r$ agente, la que vio el agente por decisión; la entropía se promedia sobre las decisiones libres.}
\label{tab:corredor_alta-ippo-vecinos-4042}
\end{table}

\begin{table}[H]
\setstretch{1}
\centering
\scriptsize
\setlength{\tabcolsep}{3pt}
\setlength{\tabcolsep}{2pt}
\begin{tabular}{rrrrrrrrrr}
\toprule
Ep. & lr & Retraso & Espera & $r$/dec. & $r$ agente & Cambios & Entropía & KL & \makecell{Eval.\\retraso} \\
 & ($\times 10^{-4}$) & (s/veh) & (s/veh) & (U2) & (U4) & de fase &  &  & (s/veh) \\
\midrule
1 & 10.0 & 207.26 & 80.12 & $-$25.94 & $-$75.89 & 674 & 0.681 & 0.0083 & -- \\
2 & 9.9 & 190.77 & 83.62 & $-$27.66 & $-$69.99 & 625 & 0.664 & 0.0090 & -- \\
3 & 9.9 & 177.18 & 72.36 & $-$24.11 & $-$64.83 & 597 & 0.616 & 0.0091 & -- \\
4 & 9.8 & 191.39 & 84.43 & $-$29.66 & $-$72.25 & 538 & 0.583 & 0.0043 & -- \\
5 & 9.8 & 183.70 & 83.30 & $-$29.80 & $-$69.50 & 489 & 0.535 & 0.0064 & -- \\
10 & 9.6 & 159.85 & 77.59 & $-$29.62 & $-$63.50 & 381 & 0.422 & 0.0070 & 167.51 \\
20 & 9.1 & 136.54 & 72.41 & $-$27.97 & $-$54.08 & 326 & 0.269 & 0.0048 & 140.45 \\
30 & 8.5 & 130.57 & 68.80 & $-$26.50 & $-$50.99 & 297 & 0.216 & 0.0064 & 137.29 \\
40 & 8.1 & 128.35 & 65.46 & $-$25.30 & $-$50.47 & 310 & 0.207 & 0.0066 & 135.65 \\
50 & 7.6 & 124.58 & 68.98 & $-$27.05 & $-$49.50 & 293 & 0.193 & 0.0066 & 128.40 \\
60 & 7.0 & 122.01 & 64.27 & $-$24.87 & $-$48.31 & 314 & 0.211 & 0.0045 & 125.58 \\
70 & 6.5 & 127.66 & 69.99 & $-$27.04 & $-$50.07 & 294 & 0.184 & 0.0041 & 129.55 \\
80 & 6.0 & 121.25 & 62.97 & $-$24.55 & $-$48.25 & 309 & 0.208 & 0.0060 & 127.82 \\
90 & 5.6 & 135.94 & 67.99 & $-$26.30 & $-$53.34 & 314 & 0.219 & 0.0036 & 127.19 \\
100 & 5.0 & 122.99 & 64.61 & $-$25.08 & $-$48.80 & 307 & 0.205 & 0.0067 & 133.38 \\
110 & 4.5 & 128.26 & 66.26 & $-$25.61 & $-$50.75 & 313 & 0.191 & 0.0039 & 121.38 \\
120 & 4.0 & 121.17 & 59.88 & $-$22.88 & $-$47.84 & 339 & 0.216 & 0.0034 & 123.41 \\
130 & 3.6 & 127.47 & 67.97 & $-$26.29 & $-$50.37 & 307 & 0.178 & 0.0029 & 123.08 \\
140 & 3.0 & 121.27 & 62.45 & $-$24.23 & $-$47.89 & 303 & 0.184 & 0.0024 & 117.96 \\
150 & 2.6 & 126.92 & 64.86 & $-$25.27 & $-$50.09 & 302 & 0.171 & 0.0024 & 121.48 \\
160 & 2.0 & 124.99 & 64.75 & $-$25.18 & $-$49.49 & 311 & 0.172 & 0.0030 & 120.60 \\
170 & 1.6 & 126.38 & 65.21 & $-$25.18 & $-$49.80 & 313 & 0.174 & 0.0025 & 119.58 \\
180 & 1.1 & 132.93 & 70.89 & $-$28.00 & $-$53.42 & 317 & 0.162 & 0.0014 & 117.90 \\
190 & 0.6 & 118.87 & 61.94 & $-$23.82 & $-$46.65 & 307 & 0.150 & 0.0017 & 121.09 \\
200 & 0.1 & 118.35 & 61.95 & $-$24.09 & $-$46.96 & 306 & 0.162 & 0.0002 & 121.44 \\
\bottomrule
\end{tabular}
\caption[Escenario corredor\_alta: entrenamiento IPPO spill, semilla base 42]{Escenario corredor\_alta: entrenamiento de IPPO spill con semilla base 42, 200 episodios (se muestran los cinco primeros y uno de cada diez; la semilla de SUMO de cada episodio es la base más su número). $r$/dec.\ es la recompensa de U2 a reloj fijo; $r$ agente, la que vio el agente por decisión; la entropía se promedia sobre las decisiones libres.}
\label{tab:corredor_alta-ippo-spill-42}
\end{table}

\begin{table}[H]
\setstretch{1}
\centering
\scriptsize
\setlength{\tabcolsep}{3pt}
\setlength{\tabcolsep}{2pt}
\begin{tabular}{rrrrrrrrrr}
\toprule
Ep. & lr & Retraso & Espera & $r$/dec. & $r$ agente & Cambios & Entropía & KL & \makecell{Eval.\\retraso} \\
 & ($\times 10^{-4}$) & (s/veh) & (s/veh) & (U2) & (U4) & de fase &  &  & (s/veh) \\
\midrule
1 & 10.0 & 204.01 & 79.96 & $-$25.63 & $-$74.05 & 681 & 0.678 & 0.0102 & -- \\
2 & 9.9 & 189.67 & 86.24 & $-$28.64 & $-$69.70 & 594 & 0.648 & 0.0098 & -- \\
3 & 9.9 & 192.05 & 88.92 & $-$30.20 & $-$70.62 & 537 & 0.593 & 0.0069 & -- \\
4 & 9.8 & 175.03 & 78.02 & $-$28.60 & $-$67.62 & 457 & 0.520 & 0.0053 & -- \\
5 & 9.8 & 176.07 & 82.96 & $-$31.09 & $-$68.47 & 439 & 0.468 & 0.0119 & -- \\
10 & 9.6 & 150.13 & 80.00 & $-$30.07 & $-$57.81 & 331 & 0.303 & 0.0065 & 152.34 \\
20 & 9.1 & 148.21 & 81.96 & $-$32.17 & $-$59.17 & 303 & 0.281 & 0.0069 & 148.35 \\
30 & 8.5 & 138.68 & 70.58 & $-$26.61 & $-$53.80 & 345 & 0.258 & 0.0089 & 136.95 \\
40 & 8.1 & 133.28 & 66.27 & $-$25.32 & $-$52.58 & 337 & 0.247 & 0.0053 & 138.69 \\
50 & 7.6 & 133.05 & 67.50 & $-$26.46 & $-$53.31 & 321 & 0.233 & 0.0078 & 132.36 \\
60 & 7.0 & 137.78 & 68.65 & $-$26.14 & $-$53.92 & 331 & 0.197 & 0.0046 & 130.19 \\
70 & 6.5 & 131.73 & 65.53 & $-$25.49 & $-$52.57 & 331 & 0.211 & 0.0046 & 125.49 \\
80 & 6.0 & 125.75 & 63.72 & $-$24.13 & $-$49.16 & 335 & 0.191 & 0.0052 & 131.35 \\
90 & 5.6 & 132.58 & 65.04 & $-$25.36 & $-$53.01 & 325 & 0.234 & 0.0061 & 122.54 \\
100 & 5.0 & 126.28 & 64.68 & $-$25.65 & $-$51.13 & 317 & 0.210 & 0.0045 & 122.52 \\
110 & 4.5 & 129.21 & 65.54 & $-$25.67 & $-$51.97 & 331 & 0.202 & 0.0034 & 131.06 \\
120 & 4.0 & 123.48 & 57.62 & $-$22.12 & $-$48.86 & 339 & 0.205 & 0.0031 & 128.76 \\
130 & 3.6 & 127.55 & 61.57 & $-$23.47 & $-$49.98 & 329 & 0.180 & 0.0033 & 129.11 \\
140 & 3.0 & 120.35 & 58.19 & $-$22.39 & $-$47.29 & 327 & 0.160 & 0.0029 & 121.67 \\
150 & 2.6 & 129.49 & 64.75 & $-$24.98 & $-$51.02 & 319 & 0.167 & 0.0020 & 126.68 \\
160 & 2.0 & 120.92 & 60.27 & $-$24.00 & $-$48.77 & 309 & 0.167 & 0.0020 & 129.54 \\
170 & 1.6 & 117.61 & 58.21 & $-$22.59 & $-$46.61 & 307 & 0.150 & 0.0018 & 120.51 \\
180 & 1.1 & 119.89 & 61.67 & $-$24.24 & $-$47.99 & 297 & 0.139 & 0.0010 & 116.44 \\
190 & 0.6 & 125.25 & 62.01 & $-$23.56 & $-$48.68 & 313 & 0.137 & 0.0006 & 122.34 \\
200 & 0.1 & 121.61 & 62.01 & $-$24.09 & $-$48.01 & 297 & 0.135 & 0.0002 & 118.69 \\
\bottomrule
\end{tabular}
\caption[Escenario corredor\_alta: entrenamiento IPPO spill, semilla base 2042]{Escenario corredor\_alta: entrenamiento de IPPO spill con semilla base 2042, 200 episodios (se muestran los cinco primeros y uno de cada diez; la semilla de SUMO de cada episodio es la base más su número). $r$/dec.\ es la recompensa de U2 a reloj fijo; $r$ agente, la que vio el agente por decisión; la entropía se promedia sobre las decisiones libres.}
\label{tab:corredor_alta-ippo-spill-2042}
\end{table}

\begin{table}[H]
\setstretch{1}
\centering
\scriptsize
\setlength{\tabcolsep}{3pt}
\setlength{\tabcolsep}{2pt}
\begin{tabular}{rrrrrrrrrr}
\toprule
Ep. & lr & Retraso & Espera & $r$/dec. & $r$ agente & Cambios & Entropía & KL & \makecell{Eval.\\retraso} \\
 & ($\times 10^{-4}$) & (s/veh) & (s/veh) & (U2) & (U4) & de fase &  &  & (s/veh) \\
\midrule
1 & 10.0 & 201.61 & 82.22 & $-$26.88 & $-$73.99 & 660 & 0.684 & 0.0058 & -- \\
2 & 9.9 & 186.96 & 77.55 & $-$25.59 & $-$68.80 & 635 & 0.661 & 0.0100 & -- \\
3 & 9.9 & 186.98 & 76.82 & $-$25.69 & $-$68.21 & 582 & 0.621 & 0.0076 & -- \\
4 & 9.8 & 173.04 & 66.94 & $-$23.08 & $-$65.49 & 572 & 0.572 & 0.0102 & -- \\
5 & 9.8 & 184.46 & 75.55 & $-$26.71 & $-$70.67 & 523 & 0.506 & 0.0126 & -- \\
10 & 9.6 & 151.60 & 71.98 & $-$27.27 & $-$59.55 & 377 & 0.391 & 0.0064 & 162.12 \\
20 & 9.1 & 147.03 & 76.71 & $-$29.68 & $-$58.65 & 343 & 0.312 & 0.0052 & 137.10 \\
30 & 8.5 & 138.92 & 67.73 & $-$25.99 & $-$55.22 & 360 & 0.281 & 0.0059 & 132.70 \\
40 & 8.1 & 143.17 & 77.29 & $-$29.61 & $-$56.47 & 322 & 0.250 & 0.0040 & 135.29 \\
50 & 7.6 & 133.84 & 73.83 & $-$28.98 & $-$53.72 & 308 & 0.218 & 0.0043 & 133.37 \\
60 & 7.0 & 133.49 & 68.30 & $-$26.46 & $-$53.06 & 324 & 0.219 & 0.0058 & 131.17 \\
70 & 6.5 & 118.51 & 58.71 & $-$22.99 & $-$47.43 & 319 & 0.193 & 0.0071 & 141.44 \\
80 & 6.0 & 136.92 & 71.61 & $-$27.77 & $-$54.27 & 314 & 0.218 & 0.0036 & 139.76 \\
90 & 5.6 & 136.19 & 71.64 & $-$28.19 & $-$54.54 & 318 & 0.212 & 0.0055 & 129.16 \\
100 & 5.0 & 128.35 & 65.68 & $-$24.98 & $-$49.74 & 307 & 0.191 & 0.0049 & 135.37 \\
110 & 4.5 & 125.97 & 65.17 & $-$24.78 & $-$48.83 & 299 & 0.177 & 0.0043 & 122.01 \\
120 & 4.0 & 127.81 & 68.57 & $-$26.46 & $-$50.06 & 291 & 0.165 & 0.0031 & 131.34 \\
130 & 3.6 & 125.26 & 65.95 & $-$25.46 & $-$49.22 & 291 & 0.155 & 0.0024 & 123.56 \\
140 & 3.0 & 120.67 & 62.00 & $-$23.89 & $-$47.25 & 305 & 0.184 & 0.0025 & 115.09 \\
150 & 2.6 & 123.82 & 65.99 & $-$25.72 & $-$49.28 & 303 & 0.177 & 0.0035 & 118.71 \\
160 & 2.0 & 120.60 & 57.59 & $-$22.12 & $-$47.40 & 322 & 0.191 & 0.0016 & 116.51 \\
170 & 1.6 & 127.39 & 66.29 & $-$25.51 & $-$49.99 & 311 & 0.163 & 0.0013 & 120.82 \\
180 & 1.1 & 122.70 & 59.99 & $-$22.94 & $-$48.07 & 322 & 0.155 & 0.0013 & 125.88 \\
190 & 0.6 & 122.73 & 60.77 & $-$23.21 & $-$48.01 & 319 & 0.160 & 0.0010 & 118.79 \\
200 & 0.1 & 124.72 & 62.09 & $-$24.31 & $-$49.50 & 307 & 0.158 & 0.0002 & 120.73 \\
\bottomrule
\end{tabular}
\caption[Escenario corredor\_alta: entrenamiento IPPO spill, semilla base 4042]{Escenario corredor\_alta: entrenamiento de IPPO spill con semilla base 4042, 200 episodios (se muestran los cinco primeros y uno de cada diez; la semilla de SUMO de cada episodio es la base más su número). $r$/dec.\ es la recompensa de U2 a reloj fijo; $r$ agente, la que vio el agente por decisión; la entropía se promedia sobre las decisiones libres.}
\label{tab:corredor_alta-ippo-spill-4042}
\end{table}



\section*{Anexo B. Código fuente}

Scripts del prototipo tal como se ejecutaron para producir los
resultados del Capítulo~\ref{cap:resultados} (SUMO 1.25.0, Python 3.14).
El código, sus comentarios y sus salidas están en español y sin tildes
en los identificadores. Se omite \texttt{graficar.py}, que solo dibuja
las figuras a partir de los mismos archivos; está en el repositorio del
proyecto.

\begin{spacing}{1}

\subsection*{rl\_semaforo.py: agente único (escenarios cruce y cruce2)}
{\fontsize{7}{8.4}\selectfont \verbatiminput{codigo/rl_semaforo.py}}

\subsection*{rl\_corredor.py: un agente independiente por semáforo (corredor y red)}
{\fontsize{7}{8.4}\selectfont \verbatiminput{codigo/rl_corredor.py}}

\subsection*{rl\_ippo.py: módulo U4, IPPO con mensajería entre vecinos}
{\fontsize{7}{8.4}\selectfont \verbatiminput{codigo/rl_ippo.py}}

\subsection*{comun.py: métricas del tripinfo, programa actuado y CSV}
{\fontsize{7}{8.4}\selectfont \verbatiminput{codigo/comun.py}}

\subsection*{evaluar.py: evaluación por semillas y estadística}
{\fontsize{7}{8.4}\selectfont \verbatiminput{codigo/evaluar.py}}

\subsection*{barrer\_fijo.py: barrido del programa de tiempo fijo}
{\fontsize{7}{8.4}\selectfont \verbatiminput{codigo/barrer_fijo.py}}

\subsection*{tablas\_tex.py: tablas del capítulo de resultados}
{\fontsize{7}{8.4}\selectfont \verbatiminput{codigo/tablas_tex.py}}

\end{spacing}
